%============================================================
%  Bd. 43 - Ringe mit Eins
%============================================================
\documentclass[main.tex]{subfiles}
\ifSubfilesClassLoaded{\BandSetupStandalone{B43}}{}
\title{Bd. 43 - Ringe mit Eins}
\author{Martin Kunze}
\date{}
\setcounter{file}{43}
\hypersetup{pdftitle={Bd. 43 - Ringe mit Eins},pdfauthor={Martin Kunze}}
\begin{document}
\providecommand{\BandXXIXProofRefStack}[1]{%
  \ensuremath{\begin{gathered}[t]#1\end{gathered}}}
\title{Bd. 43 - Ringe mit Eins}
\author{Martin Kunze}
\date{}
\hypersetup{pageanchor=false}
\maketitle
\hypersetup{pageanchor=true}
\setcounter{tocdepth}{3}
\tableofcontents
\setcounter{file}{43}
\renewcommand{\FormulaBandID}{43}

\chapter{Einleitung}

Ein Ring mit Eins verbindet eine additive abelsche Gruppe aus Band~41
mit einem multiplikativen Monoid aus Band~38 durch die beiden
Distributivgesetze. Dieser Band entwickelt Nullabsorption,
Vorzeichenregeln und Ringhomomorphismen. Anschließend werden die ganzen
Zahlen als Ring eingeordnet und ihre universelle Abbildungseigenschaft
bewiesen. Die zusätzliche Kommutativität der Multiplikation wird in
Band~44 behandelt; \(0\neq1\) wird nicht als Ringaxiom gefordert.

\chapter{Ringe mit Eins}

\section{Ringaxiome}

\FormulaAxiomDeltaK[Linksdistributivität]{%
  a,b,c\in R\vdash
  a\cdot(b+c)=a\cdot b+a\cdot c%
}{RingLeftDistributivityLaw}{%
  \DeltaRow{Mengen}{R\dsep a\dsep b\dsep c}
  \DeltaRow{Funktionen}{+\dsep\cdot}
}

\FormulaAxiomDeltaK[Rechtsdistributivität]{%
  a,b,c\in R\vdash
  (a+b)\cdot c=a\cdot c+b\cdot c%
}{RingRightDistributivityLaw}{%
  \DeltaRow{Mengen}{R\dsep a\dsep b\dsep c}
  \DeltaRow{Funktionen}{+\dsep\cdot}
}

\FormulaDefDeltaK[Begriff des Rings mit Eins]{%
  \RingAlg{R}{+}{\cdot}{0}{1}%
}{RingDef}{%
  \DeltaRow{Mengen}{R\dsep 0\dsep 1\dsep a\dsep b\dsep c}
  \DeltaRow{Funktionen}{+\dsep\cdot}
  \DeltaRow{\textbf{Axiome}}{}
  \DeltaRow{Additive abelsche Gruppe}
    {\AbelianGroupAlg{R}{+}{0}}
    [\FormulaRefAuto{AbelianGroupDef}[def]]
  \DeltaRow{Multiplikatives Monoid}
    {\MonoidAlg{R}{\cdot}{1}}
    [\FormulaRefAuto{MonoidDef}[def]]
  \DeltaRow{Linksdistributivität}{%
    \begin{aligned}[t]
      &a,b,c\in R\vdash{}\\[-2pt]
      &a\cdot(b+c)=a\cdot b+a\cdot c
    \end{aligned}}
    [\FormulaRefAuto{RingLeftDistributivityLaw}[ax]]
  \DeltaRow{Rechtsdistributivität}{%
    \begin{aligned}[t]
      &a,b,c\in R\vdash{}\\[-2pt]
      &(a+b)\cdot c=a\cdot c+b\cdot c
    \end{aligned}}
    [\FormulaRefAuto{RingRightDistributivityLaw}[ax]]
  \DeltaRow{\textbf{Neues Symbol}}{}
  \DeltaPrem{Ring mit Eins}{\RingAlg{R}{+}{\cdot}{0}{1}}
}

\FormulaAxiomDeltaK[Additiver Anteil eines Rings]{%
  \RingAlg{R}{+}{\cdot}{0}{1}
  \vdash\AbelianGroupAlg{R}{+}{0}%
}{RingAdditiveAbelianGroupAxiom}{%
  \DeltaRow{Mengen}{R\dsep 0\dsep 1}
  \DeltaRow{Funktionen}{+\dsep\cdot}
}

\FormulaAxiomDeltaK[Multiplikativer Anteil eines Rings]{%
  \RingAlg{R}{+}{\cdot}{0}{1}
  \vdash\MonoidAlg{R}{\cdot}{1}%
}{RingMultiplicativeMonoidAxiom}{%
  \DeltaRow{Mengen}{R\dsep 0\dsep 1}
  \DeltaRow{Funktionen}{+\dsep\cdot}
}

\FormulaAxiomDeltaK[Linkes Distributivgesetz eines Rings]{%
  \RingAlg{R}{+}{\cdot}{0}{1}\dsep a,b,c\in R
  \vdash a\cdot(b+c)=a\cdot b+a\cdot c%
}{RingLeftDistributivityAxiom}{%
  \DeltaRow{Mengen}{R\dsep 0\dsep 1\dsep a\dsep b\dsep c}
  \DeltaRow{Funktionen}{+\dsep\cdot}
}

\FormulaAxiomDeltaK[Rechtes Distributivgesetz eines Rings]{%
  \RingAlg{R}{+}{\cdot}{0}{1}\dsep a,b,c\in R
  \vdash (a+b)\cdot c=a\cdot c+b\cdot c%
}{RingRightDistributivityAxiom}{%
  \DeltaRow{Mengen}{R\dsep 0\dsep 1\dsep a\dsep b\dsep c}
  \DeltaRow{Funktionen}{+\dsep\cdot}
}

\section{Abgeleitete Trägergesetze}

\FormulaThmDeltaK[Additive Gruppe eines Rings]{%
  \RingAlg{R}{+}{\cdot}{0}{1}\vdash\GroupAlg{R}{+}{0}%
}{RingAdditiveGroup}{%
  \DeltaRow{Mengen}{R\dsep 0\dsep 1}
  \DeltaRow{Funktionen}{+\dsep\cdot}
}
\begin{tabproof}
  \proofstep{1}{\RingAlg{R}{+}{\cdot}{0}{1}}{\rA}
  \proofstep{1}{\AbelianGroupAlg{R}{+}{0}}{%
    \FormulaRefAuto{RingAdditiveAbelianGroupAxiom}[ax]{1}}
  \proofstep{1}{\GroupAlg{R}{+}{0}}{%
    \FormulaRefAuto{AbelianGroupBaseGroupAxiom}[ax]{2}}
\end{tabproof}

\FormulaThmDeltaK[Multiplikative Halbgruppe eines Rings]{%
  \RingAlg{R}{+}{\cdot}{0}{1}\vdash\SemigroupAlg{R}{\cdot}%
}{RingMultiplicativeSemigroup}{%
  \DeltaRow{Mengen}{R\dsep 0\dsep 1}
  \DeltaRow{Funktionen}{+\dsep\cdot}
}
\begin{tabproof}
  \proofstep{1}{\RingAlg{R}{+}{\cdot}{0}{1}}{\rA}
  \proofstep{1}{\MonoidAlg{R}{\cdot}{1}}{%
    \FormulaRefAuto{RingMultiplicativeMonoidAxiom}[ax]{1}}
  \proofstep{1}{\SemigroupAlg{R}{\cdot}}{%
    \FormulaRefAuto{MonoidBaseSemigroupAxiom}[ax]{2}}
\end{tabproof}

\FormulaThmDeltaK[Abgeschlossenheit der Ringmultiplikation]{%
  \RingAlg{R}{+}{\cdot}{0}{1}\dsep a,b\in R
  \vdash a\cdot b\in R%
}{RingMultiplicationClosure}{%
  \DeltaRow{Mengen}{R\dsep 0\dsep 1\dsep a\dsep b}
  \DeltaRow{Funktionen}{+\dsep\cdot}
}
\begin{tabproof}
  \proofstep{1}{\RingAlg{R}{+}{\cdot}{0}{1}}{\rA}
  \proofstep{2}{a\in R}{\rA}
  \proofstep{3}{b\in R}{\rA}
  \proofstep{1}{\SemigroupAlg{R}{\cdot}}{%
    \FormulaRefAuto{RingMultiplicativeSemigroup}[thm]{1}}
  \proofstep{1,2,3}{a\cdot b\in R}{%
    \FormulaRefAuto{SemigroupClosureAxiom}[ax]{4,2,3}}
\end{tabproof}

\section{Additives Inverses und Vorzeichen}

\FormulaDefDeltaK[Additives Inverses in einem Ring]{%
  \begin{aligned}[t]
    &\RingAlg{R}{+}{\cdot}{0}{1}\dsep a\in R\vdash{}\\
    &-a\coloneqq
      \iota b\,
      \bigl(b\in R\land a+b=0\land b+a=0\bigr)
  \end{aligned}
}{RingNegativeElement}{%
  \DeltaRow{Mengen}{R\dsep 0\dsep 1\dsep a\dsep b}
  \DeltaRow{Funktionen}{+\dsep\cdot}
  \DeltaRow{Träger}{-a\in R}
  \DeltaRow{Rechtsinverse}{a+(-a)=0}
  \DeltaRow{Linksinverse}{(-a)+a=0}
  \DeltaRow{Eindeutigkeit}{\FormulaRefAuto{GroupInverseUnique}[thm]}
  \DeltaRow{\textbf{Neues Symbol}}{}
  \DeltaPrem{Additives Inverses}{-a}
}

\FormulaThmDeltaK[Additives Inverses liegt im Ringträger]{%
  \RingAlg{R}{+}{\cdot}{0}{1}\dsep a\in R\vdash -a\in R%
}{RingNegativeCarrier}{%
  \DeltaRow{Mengen}{R\dsep 0\dsep 1\dsep a}
  \DeltaRow{Funktionen}{+\dsep\cdot}
}
\begin{tabproof}
  \proofstep{1}{\RingAlg{R}{+}{\cdot}{0}{1}}{\rA}
  \proofstep{2}{a\in R}{\rA}
  \proofstep{1,2}{-a\in R}{%
    \FormulaRefAuto{RingNegativeElement}[def]{1,2}}
\end{tabproof}

\FormulaThmDeltaK[Rechte additive Inversengleichung]{%
  \RingAlg{R}{+}{\cdot}{0}{1}\dsep a\in R
  \vdash a+(-a)=0%
}{RingNegativeRight}{%
  \DeltaRow{Mengen}{R\dsep 0\dsep 1\dsep a}
  \DeltaRow{Funktionen}{+\dsep\cdot}
}
\begin{tabproof}
  \proofstep{1}{\RingAlg{R}{+}{\cdot}{0}{1}}{\rA}
  \proofstep{2}{a\in R}{\rA}
  \proofstep{1,2}{a+(-a)=0}{%
    \FormulaRefAuto{RingNegativeElement}[def]{1,2}}
\end{tabproof}

\FormulaThmDeltaK[Linke additive Inversengleichung]{%
  \RingAlg{R}{+}{\cdot}{0}{1}\dsep a\in R
  \vdash (-a)+a=0%
}{RingNegativeLeft}{%
  \DeltaRow{Mengen}{R\dsep 0\dsep 1\dsep a}
  \DeltaRow{Funktionen}{+\dsep\cdot}
}
\begin{tabproof}
  \proofstep{1}{\RingAlg{R}{+}{\cdot}{0}{1}}{\rA}
  \proofstep{2}{a\in R}{\rA}
  \proofstep{1,2}{(-a)+a=0}{%
    \FormulaRefAuto{RingNegativeElement}[def]{1,2}}
\end{tabproof}

\FormulaThmDeltaK[Doppelte Negation]{%
  \RingAlg{R}{+}{\cdot}{0}{1}\dsep a\in R
  \vdash -(-a)=a%
}{RingDoubleNegative}{%
  \DeltaRow{Mengen}{R\dsep 0\dsep 1\dsep a}
  \DeltaRow{Funktionen}{+\dsep\cdot}
}
\begin{tabproofwide}
  \proofstepwidestar[1]{\RingAlg{R}{+}{\cdot}{0}{1}}{\rA}
  \proofstepwidestar[2]{a\in R}{\rA}
  \proofstepwidestar[1]{\GroupAlg{R}{+}{0}}{%
    \FormulaRefAuto{RingAdditiveGroup}[thm]{1}}
  \proofstepwidestar[1,2]{-a\in R}{%
    \FormulaRefAuto{RingNegativeCarrier}[thm]{1,2}}
  \proofstepwidestar[1,4]{-(-a)\in R}{%
    \FormulaRefAuto{RingNegativeCarrier}[thm]{1,4}}
  \proofstepwidestar[1,4]{-(-a)+(-a)=0}{%
    \FormulaRefAuto{RingNegativeLeft}[thm]{1,4}}
  \proofstepwidestar[1,2]{(-a)+a=0}{%
    \FormulaRefAuto{RingNegativeLeft}[thm]{1,2}}
  \proofstepwidestar[1,2,3,4,5,6,7]{-(-a)=a}{%
    \FormulaRefAuto{GroupInverseUnique}[thm]{3,4,5,2,6,7}}
\end{tabproofwide}

\section{Nullabsorption}

\FormulaThmDeltaK[Linke Nullabsorption]{%
  \RingAlg{R}{+}{\cdot}{0}{1}\dsep a\in R
  \vdash 0\cdot a=0%
}{RingLeftZeroAbsorption}{%
  \DeltaRow{Mengen}{R\dsep 0\dsep 1\dsep a}
  \DeltaRow{Funktionen}{+\dsep\cdot}
}
\begin{tabproofsplitwide}
  \proofpartwideindR[Verdopplung]{%
    \RingAlg{R}{+}{\cdot}{0}{1}\dsep a\in R
    \vdash 0\cdot a=0\cdot a+0\cdot a
  }{%
    \RingAlg{R}{+}{\cdot}{0}{1}\dsep a\in R
    \vdash 0\cdot a=0\cdot a+0\cdot a%
  }[RingLeftZeroDoubling]
    \proofstepwidestar[1]{\RingAlg{R}{+}{\cdot}{0}{1}}{\rA}
    \proofstepwidestar[2]{a\in R}{\rA}
    \proofstepwidestar[1]{\GroupAlg{R}{+}{0}}{%
      \FormulaRefAuto{RingAdditiveGroup}[thm]{1}}
    \proofstepwidestar[1]{0\in R}{%
      \FormulaRefAuto{GroupIdentityCarrier}[thm]{3}}
    \proofstepwide[3,4]{0+0}{=}{0}{%
      \FormulaRefAuto{GroupRightIdentity}[thm]{3,4}}
    \proofstepwide[5]{0\cdot a}{=}{(0+0)\cdot a}{%
      \rIE{\FormulaRefAuto{a=b\vdash b=a}{5}}}
    \proofstepwide[1,2,4]{(0+0)\cdot a}{=}{%
      0\cdot a+0\cdot a}{%
      \FormulaRefAuto{RingRightDistributivityAxiom}[ax]{1,4,4,2}}
    \proofstepwidestar[1,2]{0\cdot a=0\cdot a+0\cdot a}{%
      \rChain{6,7}}
  \closeproofpartwideindR

  \proofpartwide{Kürzung}
    \proofstepwidestar[1]{\RingAlg{R}{+}{\cdot}{0}{1}}{\rA}
    \proofstepwidestar[2]{a\in R}{\rA}
    \proofstepwidestar[1]{\GroupAlg{R}{+}{0}}{%
      \FormulaRefAuto{RingAdditiveGroup}[thm]{1}}
    \proofstepwidestar[1,2]{0\in R}{%
      \FormulaRefAuto{GroupIdentityCarrier}[thm]{3}}
    \proofstepwidestar[1,2]{0\cdot a\in R}{%
      \FormulaRefAuto{RingMultiplicationClosure}[thm]{%
        1,4,2}}
    \proofstepwide[3,5]{0\cdot a+0}{=}{0\cdot a}{%
      \FormulaRefAuto{GroupRightIdentity}[thm]{3,5}}
    \proofstepwide[1,2]{0\cdot a}{=}{0\cdot a+0\cdot a}{%
      \FormulaRefAuto{RingLeftZeroDoubling}[thm]{1,2}}
    \proofstepwidestar[1,2,3,5,6,7]{%
      0\cdot a+0=0\cdot a+0\cdot a}{\rChain{6,7}}
    \proofstepwidestar[1,2,3,5,8]{0\in R}{%
      \FormulaRefAuto{GroupIdentityCarrier}[thm]{3}}
    \proofstepwidestar[1,2,3,5,8]{0=0\cdot a}{%
      \FormulaRefAuto{GroupLeftCancellation}[thm]{3,5,
        9,5,8}}
    \proofstepwidestar[1,2]{0\cdot a=0}{%
      \FormulaRefAuto{a=b\vdash b=a}{10}}
  \closeproofpartwide
\end{tabproofsplitwide}

\FormulaThmDeltaK[Rechte Nullabsorption]{%
  \RingAlg{R}{+}{\cdot}{0}{1}\dsep a\in R
  \vdash a\cdot0=0%
}{RingRightZeroAbsorption}{%
  \DeltaRow{Mengen}{R\dsep 0\dsep 1\dsep a}
  \DeltaRow{Funktionen}{+\dsep\cdot}
}
\begin{tabproofsplitwide}
  \proofpartwideindR[Verdopplung]{%
    \RingAlg{R}{+}{\cdot}{0}{1}\dsep a\in R
    \vdash a\cdot0=a\cdot0+a\cdot0
  }{%
    \RingAlg{R}{+}{\cdot}{0}{1}\dsep a\in R
    \vdash a\cdot0=a\cdot0+a\cdot0%
  }[RingRightZeroDoubling]
    \proofstepwidestar[1]{\RingAlg{R}{+}{\cdot}{0}{1}}{\rA}
    \proofstepwidestar[2]{a\in R}{\rA}
    \proofstepwidestar[1]{\GroupAlg{R}{+}{0}}{%
      \FormulaRefAuto{RingAdditiveGroup}[thm]{1}}
    \proofstepwidestar[1]{0\in R}{%
      \FormulaRefAuto{GroupIdentityCarrier}[thm]{3}}
    \proofstepwide[3,4]{0+0}{=}{0}{%
      \FormulaRefAuto{GroupRightIdentity}[thm]{3,4}}
    \proofstepwide[5]{a\cdot0}{=}{a\cdot(0+0)}{%
      \rIE{\FormulaRefAuto{a=b\vdash b=a}{5}}}
    \proofstepwide[1,2,4]{a\cdot(0+0)}{=}{%
      a\cdot0+a\cdot0}{%
      \FormulaRefAuto{RingLeftDistributivityAxiom}[ax]{1,2,4,4}}
    \proofstepwidestar[1,2]{a\cdot0=a\cdot0+a\cdot0}{%
      \rChain{6,7}}
  \closeproofpartwideindR

  \proofpartwide{Kürzung}
    \proofstepwidestar[1]{\RingAlg{R}{+}{\cdot}{0}{1}}{\rA}
    \proofstepwidestar[2]{a\in R}{\rA}
    \proofstepwidestar[1]{\GroupAlg{R}{+}{0}}{%
      \FormulaRefAuto{RingAdditiveGroup}[thm]{1}}
    \proofstepwidestar[1,2]{0\in R}{%
      \FormulaRefAuto{GroupIdentityCarrier}[thm]{3}}
    \proofstepwidestar[1,2]{a\cdot0\in R}{%
      \FormulaRefAuto{RingMultiplicationClosure}[thm]{%
        1,2,4}}
    \proofstepwide[3,5]{a\cdot0+0}{=}{a\cdot0}{%
      \FormulaRefAuto{GroupRightIdentity}[thm]{3,5}}
    \proofstepwide[1,2]{a\cdot0}{=}{a\cdot0+a\cdot0}{%
      \FormulaRefAuto{RingRightZeroDoubling}[thm]{1,2}}
    \proofstepwidestar[1,2,3,5,6,7]{%
      a\cdot0+0=a\cdot0+a\cdot0}{\rChain{6,7}}
    \proofstepwidestar[1,2,3,5,8]{0\in R}{%
      \FormulaRefAuto{GroupIdentityCarrier}[thm]{3}}
    \proofstepwidestar[1,2,3,5,8]{0=a\cdot0}{%
      \FormulaRefAuto{GroupLeftCancellation}[thm]{3,5,
        9,5,8}}
    \proofstepwidestar[1,2]{a\cdot0=0}{%
      \FormulaRefAuto{a=b\vdash b=a}{10}}
  \closeproofpartwide
\end{tabproofsplitwide}

\section{Vorzeichenregeln}

\FormulaThmDeltaK[Negatives Vorzeichen im linken Faktor]{%
  \RingAlg{R}{+}{\cdot}{0}{1}\dsep a,b\in R
  \vdash (-a)\cdot b=-(a\cdot b)%
}{RingLeftNegativeProduct}{%
  \DeltaRow{Mengen}{R\dsep 0\dsep 1\dsep a\dsep b}
  \DeltaRow{Funktionen}{+\dsep\cdot}
}
\begin{tabproofsplitwide}
  \proofpartwideindR[Rechtsinverse Eigenschaft]{%
    \begin{aligned}[t]
      &\RingAlg{R}{+}{\cdot}{0}{1}\dsep a,b\in R\vdash{}\\
      &a\cdot b+(-a)\cdot b=0
    \end{aligned}
  }{%
    \RingAlg{R}{+}{\cdot}{0}{1}\dsep a,b\in R
    \vdash a\cdot b+(-a)\cdot b=0%
  }[RingLeftNegativeRightInverse]
    \proofstepwidestar[1]{\RingAlg{R}{+}{\cdot}{0}{1}}{\rA}
    \proofstepwidestar[2]{a\in R}{\rA}
    \proofstepwidestar[3]{b\in R}{\rA}
    \proofstepwidestar[1,2]{-a\in R}{%
      \FormulaRefAuto{RingNegativeCarrier}[thm]{1,2}}
    \proofstepwidestar[1,2,3,4]{(a+(-a))\cdot b=a\cdot b+(-a)\cdot b}{%
      \FormulaRefAuto{RingRightDistributivityAxiom}[ax]{1,2,4,3}}
    \proofstepwide[1,2,3,4]{a\cdot b+(-a)\cdot b}{=}{%
      (a+(-a))\cdot b}{%
      \FormulaRefAuto{a=b\vdash b=a}{%
        5}}
    \proofstepwide[1,2,3]{(a+(-a))\cdot b}{=}{0\cdot b}{%
      \FormulaRefAuto{RingNegativeRight}[thm]{1,2}}
    \proofstepwide[1,3]{0\cdot b}{=}{0}{%
      \FormulaRefAuto{RingLeftZeroAbsorption}[thm]{1,3}}
    \proofstepwidestar[1,2,3]{a\cdot b+(-a)\cdot b=0}{%
      \rChain{6,7,8}}
  \closeproofpartwideindR

  \proofpartwideindR[Linksinverse Eigenschaft]{%
    \begin{aligned}[t]
      &\RingAlg{R}{+}{\cdot}{0}{1}\dsep a,b\in R\vdash{}\\
      &(-a)\cdot b+a\cdot b=0
    \end{aligned}
  }{%
    \RingAlg{R}{+}{\cdot}{0}{1}\dsep a,b\in R
    \vdash (-a)\cdot b+a\cdot b=0%
  }[RingLeftNegativeLeftInverse]
    \proofstepwidestar[1]{\RingAlg{R}{+}{\cdot}{0}{1}}{\rA}
    \proofstepwidestar[2]{a\in R}{\rA}
    \proofstepwidestar[3]{b\in R}{\rA}
    \proofstepwidestar[1,2]{-a\in R}{%
      \FormulaRefAuto{RingNegativeCarrier}[thm]{1,2}}
    \proofstepwidestar[1,2,3,4]{((-a)+a)\cdot b=(-a)\cdot b+a\cdot b}{%
      \FormulaRefAuto{RingRightDistributivityAxiom}[ax]{1,4,2,3}}
    \proofstepwide[1,2,3,4]{(-a)\cdot b+a\cdot b}{=}{%
      ((-a)+a)\cdot b}{%
      \FormulaRefAuto{a=b\vdash b=a}{%
        5}}
    \proofstepwide[1,2,3]{((-a)+a)\cdot b}{=}{0\cdot b}{%
      \FormulaRefAuto{RingNegativeLeft}[thm]{1,2}}
    \proofstepwide[1,3]{0\cdot b}{=}{0}{%
      \FormulaRefAuto{RingLeftZeroAbsorption}[thm]{1,3}}
    \proofstepwidestar[1,2,3]{(-a)\cdot b+a\cdot b=0}{%
      \rChain{6,7,8}}
  \closeproofpartwideindR

  \proofpartwide{Eindeutigkeit}
    \proofstepwidestar[1]{\RingAlg{R}{+}{\cdot}{0}{1}}{\rA}
    \proofstepwidestar[2]{a\in R}{\rA}
    \proofstepwidestar[3]{b\in R}{\rA}
    \proofstepwidestar[1]{\GroupAlg{R}{+}{0}}{%
      \FormulaRefAuto{RingAdditiveGroup}[thm]{1}}
    \proofstepwidestar[1,2,3]{a\cdot b\in R}{%
      \FormulaRefAuto{RingMultiplicationClosure}[thm]{1,2,3}}
    \proofstepwidestar[1,2,3]{-a\in R}{%
      \FormulaRefAuto{RingNegativeCarrier}[thm]{1,2}}
    \proofstepwidestar[1,2,3]{(-a)\cdot b\in R}{%
      \FormulaRefAuto{RingMultiplicationClosure}[thm]{%
        1,6,3}}
    \proofstepwidestar[1,5]{-(a\cdot b)\in R}{%
      \FormulaRefAuto{RingNegativeCarrier}[thm]{1,5}}
    \proofstepwidestar[1,2,3]{(-a)\cdot b+a\cdot b=0}{%
      \FormulaRefAuto{RingLeftNegativeLeftInverse}[thm]{1,2,3}}
    \proofstepwidestar[1,5]{a\cdot b+(-(a\cdot b))=0}{%
      \FormulaRefAuto{RingNegativeRight}[thm]{1,5}}
    \proofstepwidestar[1,2,3,4,5,7,8,9,10]{%
      (-a)\cdot b=-(a\cdot b)}{%
      \FormulaRefAuto{GroupInverseUnique}[thm]{4,5,7,8,9,10}}
  \closeproofpartwide
\end{tabproofsplitwide}

\FormulaThmDeltaK[Negatives Vorzeichen im rechten Faktor]{%
  \RingAlg{R}{+}{\cdot}{0}{1}\dsep a,b\in R
  \vdash a\cdot(-b)=-(a\cdot b)%
}{RingRightNegativeProduct}{%
  \DeltaRow{Mengen}{R\dsep 0\dsep 1\dsep a\dsep b}
  \DeltaRow{Funktionen}{+\dsep\cdot}
}
\begin{tabproofsplitwide}
  \proofpartwideindR[Rechtsinverse Eigenschaft]{%
    \begin{aligned}[t]
      &\RingAlg{R}{+}{\cdot}{0}{1}\dsep a,b\in R\vdash{}\\
      &a\cdot b+a\cdot(-b)=0
    \end{aligned}
  }{%
    \RingAlg{R}{+}{\cdot}{0}{1}\dsep a,b\in R
    \vdash a\cdot b+a\cdot(-b)=0%
  }[RingRightNegativeRightInverse]
    \proofstepwidestar[1]{\RingAlg{R}{+}{\cdot}{0}{1}}{\rA}
    \proofstepwidestar[2]{a\in R}{\rA}
    \proofstepwidestar[3]{b\in R}{\rA}
    \proofstepwidestar[1,3]{-b\in R}{%
      \FormulaRefAuto{RingNegativeCarrier}[thm]{1,3}}
    \proofstepwidestar[1,2,3,4]{a\cdot(b+(-b))=a\cdot b+a\cdot(-b)}{%
      \FormulaRefAuto{RingLeftDistributivityAxiom}[ax]{1,2,3,4}}
    \proofstepwide[1,2,3,4]{a\cdot b+a\cdot(-b)}{=}{%
      a\cdot(b+(-b))}{%
      \FormulaRefAuto{a=b\vdash b=a}{%
        5}}
    \proofstepwide[1,2,3]{a\cdot(b+(-b))}{=}{a\cdot0}{%
      \FormulaRefAuto{RingNegativeRight}[thm]{1,3}}
    \proofstepwide[1,2]{a\cdot0}{=}{0}{%
      \FormulaRefAuto{RingRightZeroAbsorption}[thm]{1,2}}
    \proofstepwidestar[1,2,3]{a\cdot b+a\cdot(-b)=0}{%
      \rChain{6,7,8}}
  \closeproofpartwideindR

  \proofpartwideindR[Linksinverse Eigenschaft]{%
    \begin{aligned}[t]
      &\RingAlg{R}{+}{\cdot}{0}{1}\dsep a,b\in R\vdash{}\\
      &a\cdot(-b)+a\cdot b=0
    \end{aligned}
  }{%
    \RingAlg{R}{+}{\cdot}{0}{1}\dsep a,b\in R
    \vdash a\cdot(-b)+a\cdot b=0%
  }[RingRightNegativeLeftInverse]
    \proofstepwidestar[1]{\RingAlg{R}{+}{\cdot}{0}{1}}{\rA}
    \proofstepwidestar[2]{a\in R}{\rA}
    \proofstepwidestar[3]{b\in R}{\rA}
    \proofstepwidestar[1,3]{-b\in R}{%
      \FormulaRefAuto{RingNegativeCarrier}[thm]{1,3}}
    \proofstepwidestar[1,2,3,4]{a\cdot((-b)+b)=a\cdot(-b)+a\cdot b}{%
      \FormulaRefAuto{RingLeftDistributivityAxiom}[ax]{1,2,4,3}}
    \proofstepwide[1,2,3,4]{a\cdot(-b)+a\cdot b}{=}{%
      a\cdot((-b)+b)}{%
      \FormulaRefAuto{a=b\vdash b=a}{%
        5}}
    \proofstepwide[1,2,3]{a\cdot((-b)+b)}{=}{a\cdot0}{%
      \FormulaRefAuto{RingNegativeLeft}[thm]{1,3}}
    \proofstepwide[1,2]{a\cdot0}{=}{0}{%
      \FormulaRefAuto{RingRightZeroAbsorption}[thm]{1,2}}
    \proofstepwidestar[1,2,3]{a\cdot(-b)+a\cdot b=0}{%
      \rChain{6,7,8}}
  \closeproofpartwideindR

  \proofpartwide{Eindeutigkeit}
    \proofstepwidestar[1]{\RingAlg{R}{+}{\cdot}{0}{1}}{\rA}
    \proofstepwidestar[2]{a\in R}{\rA}
    \proofstepwidestar[3]{b\in R}{\rA}
    \proofstepwidestar[1]{\GroupAlg{R}{+}{0}}{%
      \FormulaRefAuto{RingAdditiveGroup}[thm]{1}}
    \proofstepwidestar[1,2,3]{a\cdot b\in R}{%
      \FormulaRefAuto{RingMultiplicationClosure}[thm]{1,2,3}}
    \proofstepwidestar[1,2,3]{-b\in R}{%
      \FormulaRefAuto{RingNegativeCarrier}[thm]{1,3}}
    \proofstepwidestar[1,2,3]{a\cdot(-b)\in R}{%
      \FormulaRefAuto{RingMultiplicationClosure}[thm]{%
        1,2,6}}
    \proofstepwidestar[1,5]{-(a\cdot b)\in R}{%
      \FormulaRefAuto{RingNegativeCarrier}[thm]{1,5}}
    \proofstepwidestar[1,2,3]{a\cdot(-b)+a\cdot b=0}{%
      \FormulaRefAuto{RingRightNegativeLeftInverse}[thm]{1,2,3}}
    \proofstepwidestar[1,5]{a\cdot b+(-(a\cdot b))=0}{%
      \FormulaRefAuto{RingNegativeRight}[thm]{1,5}}
    \proofstepwidestar[1,2,3,4,5,7,8,9,10]{%
      a\cdot(-b)=-(a\cdot b)}{%
      \FormulaRefAuto{GroupInverseUnique}[thm]{4,5,7,8,9,10}}
  \closeproofpartwide
\end{tabproofsplitwide}

\FormulaThmDeltaK[Produkt zweier negativer Faktoren]{%
  \RingAlg{R}{+}{\cdot}{0}{1}\dsep a,b\in R
  \vdash (-a)\cdot(-b)=a\cdot b%
}{RingNegativeTimesNegative}{%
  \DeltaRow{Mengen}{R\dsep 0\dsep 1\dsep a\dsep b}
  \DeltaRow{Funktionen}{+\dsep\cdot}
}
\begin{tabproofwide}
  \proofstepwidestar[1]{\RingAlg{R}{+}{\cdot}{0}{1}}{\rA}
  \proofstepwidestar[2]{a\in R}{\rA}
  \proofstepwidestar[3]{b\in R}{\rA}
  \proofstepwidestar[1,2,3]{-b\in R}{%
      \FormulaRefAuto{RingNegativeCarrier}[thm]{1,3}}
  \proofstepwide[1,2,3]{(-a)\cdot(-b)}{=}{-(a\cdot(-b))}{%
    \FormulaRefAuto{RingLeftNegativeProduct}[thm]{%
      1,2,4}}
  \proofstepwide[1,2,3]{-(a\cdot(-b))}{=}{-(-(a\cdot b))}{%
    \FormulaRefAuto{RingRightNegativeProduct}[thm]{1,2,3}}
  \proofstepwidestar[1,2,3]{a\cdot b\in R}{%
      \FormulaRefAuto{RingMultiplicationClosure}[thm]{1,2,3}}
  \proofstepwide[1,2,3]{-(-(a\cdot b))}{=}{a\cdot b}{%
    \FormulaRefAuto{RingDoubleNegative}[thm]{%
      1,7}}
  \proofstepwidestar[1,2,3]{(-a)\cdot(-b)=a\cdot b}{%
    \rChain{5,6,8}}
\end{tabproofwide}

\FormulaThmDeltaK[Additives Inverses einer Summe]{%
  \begin{aligned}[t]&\RingAlg{R}{\oplus}{\odot}{0_R}{1_R}\dsep p,q\in R\vdash{}\\&-(p\oplus q)=(-p)\oplus(-q)\end{aligned}
}{RingNegativeSum}{%
  \DeltaRow{Mengen}{R\dsep0_R\dsep1_R\dsep p\dsep q}
  \DeltaRow{Funktionen}{\oplus\dsep\odot}
}
\begingroup
\ProofWideReasonsNearFormula
\begin{tabproofsplitwide}
  \proofpartwide{Träger}
  \proofstepwidestarbreakkeep[1]{\RingAlg{R}{\oplus}{\odot}{0_R}{1_R}}{%
    \rA}
  \proofstepwidestarbreakkeep[2]{p\in R}{%
    \rA}
  \proofstepwidestarbreakkeep[3]{q\in R}{%
    \rA}
  \proofstepwidestarbreakkeep[1]{\AbelianGroupAlg{R}{\oplus}{0_R}}{%
    \FormulaRefAuto{RingAdditiveAbelianGroupAxiom}[ax]{1}}
  \proofstepwidestarbreakkeep[1]{\GroupAlg{R}{\oplus}{0_R}}{%
    \FormulaRefAuto{RingAdditiveGroup}[thm]{1}}
  \proofstepwidestarbreakkeep[1]{\CommutativeMonoidAlg{R}{\oplus}{0_R}}{%
    \FormulaRefAuto{AbelianGroupIsCommutativeMonoid}[thm]{4}}
  \proofstepwidestarbreakkeep[1]{\CommutativeSemigroupAlg{R}{\oplus}}{%
    \FormulaRefAuto{CommutativeMonoidIsCommutativeSemigroup}[thm]{6}}
  \proofstepwidestarbreakkeep[1,2]{-p\in R}{%
    \FormulaRefAuto{RingNegativeCarrier}[thm]{1,2}}
  \proofstepwidestarbreakkeep[1,3]{-q\in R}{%
    \FormulaRefAuto{RingNegativeCarrier}[thm]{1,3}}
  \proofstepwidestarbreakkeep[1,2,3]{p\oplus q\in R}{%
    \FormulaRefAuto{GroupClosure}[thm]{5,2,3}}
  \proofstepwidestarbreakkeep[1,2,3]{-(p\oplus q)\in R}{%
    \FormulaRefAuto{RingNegativeCarrier}[thm]{1,10}}
  \proofstepwidestarbreakkeep[1,2,3]{(-p)\oplus(-q)\in R}{%
    \FormulaRefAuto{GroupClosure}[thm]{5,8,9}}
  \closeproofpartwide

  \proofpartwide{Zwei inverse Elemente derselben Summe}
  \setcounter{proofstepnr}{12}
  \proofstepwidestarbreakkeep[1,2,3]{(p\oplus q)\oplus((-p)\oplus(-q))=(p\oplus(-p))\oplus(q\oplus(-q))}{%
    \FormulaRefAuto{CommutativeSemigroupFourTermExchange}[thm]{7,2,3,8,9}}
  \proofstepwidestarbreakkeep[1,2]{p\oplus(-p)=0_R}{%
    \FormulaRefAuto{RingNegativeRight}[thm]{1,2}}
  \proofstepwidestarbreakkeep[1,3]{q\oplus(-q)=0_R}{%
    \FormulaRefAuto{RingNegativeRight}[thm]{1,3}}
  \proofstepwidestarbreakkeep[1]{0_R\in R}{%
    \FormulaRefAuto{GroupIdentityCarrier}[thm]{5}}
  \proofstepwidestarbreakkeep[1]{0_R\oplus0_R=0_R}{%
    \FormulaRefAuto{GroupLeftIdentity}[thm]{5,16}}
  \proofstepwidestarbreakkeep[1,2,3]{(p\oplus q)\oplus((-p)\oplus(-q))=0_R}{%
    \rIE{14,15,17,13}}
  \proofstepwidestarbreakkeep[1,2,3]{(-(p\oplus q))\oplus(p\oplus q)=0_R}{%
    \FormulaRefAuto{RingNegativeLeft}[thm]{1,10}}
  \proofstepwidestarbreakkeep[1,2,3]{-(p\oplus q)=(-p)\oplus(-q)}{%
    \FormulaRefAuto{GroupInverseUnique}[thm]{5,10,11,12,19,18}}
  \closeproofpartwide
\end{tabproofsplitwide}
\endgroup

\FormulaThmDeltaK[Addition zweier Differenzen]{%
  \begin{aligned}[t]&\RingAlg{R}{\oplus}{\odot}{0_R}{1_R}\dsep p,q,r,s\in R\vdash{}\\&(p\oplus(-q))\oplus(r\oplus(-s))=(p\oplus r)\oplus(-(q\oplus s))\end{aligned}
}{RingDifferenceAddition}{%
  \DeltaRow{Mengen}{R\dsep0_R\dsep1_R\dsep p\dsep q\dsep r\dsep s}
  \DeltaRow{Funktionen}{\oplus\dsep\odot}
}
\begingroup
\ProofWideReasonsNearFormula
\begin{tabproofwide}
  \proofstepwidestarbreakkeep[1]{\RingAlg{R}{\oplus}{\odot}{0_R}{1_R}}{%
    \rA}
  \proofstepwidestarbreakkeep[2]{p\in R}{%
    \rA}
  \proofstepwidestarbreakkeep[3]{q\in R}{%
    \rA}
  \proofstepwidestarbreakkeep[4]{r\in R}{%
    \rA}
  \proofstepwidestarbreakkeep[5]{s\in R}{%
    \rA}
  \proofstepwidestarbreakkeep[1]{\AbelianGroupAlg{R}{\oplus}{0_R}}{%
    \FormulaRefAuto{RingAdditiveAbelianGroupAxiom}[ax]{1}}
  \proofstepwidestarbreakkeep[1]{\GroupAlg{R}{\oplus}{0_R}}{%
    \FormulaRefAuto{RingAdditiveGroup}[thm]{1}}
  \proofstepwidestarbreakkeep[1]{\CommutativeMonoidAlg{R}{\oplus}{0_R}}{%
    \FormulaRefAuto{AbelianGroupIsCommutativeMonoid}[thm]{6}}
  \proofstepwidestarbreakkeep[1]{\CommutativeSemigroupAlg{R}{\oplus}}{%
    \FormulaRefAuto{CommutativeMonoidIsCommutativeSemigroup}[thm]{8}}
  \proofstepwidestarbreakkeep[1,3]{-q\in R}{%
    \FormulaRefAuto{RingNegativeCarrier}[thm]{1,3}}
  \proofstepwidestarbreakkeep[1,5]{-s\in R}{%
    \FormulaRefAuto{RingNegativeCarrier}[thm]{1,5}}
  \proofstepwidestarbreakkeep[1,2,3,4,5]{(p\oplus(-q))\oplus(r\oplus(-s))=(p\oplus r)\oplus((-q)\oplus(-s))}{%
    \FormulaRefAuto{CommutativeSemigroupFourTermExchange}[thm]{9,2,10,4,11}}
  \proofstepwidestarbreakkeep[1,3,5]{-(q\oplus s)=(-q)\oplus(-s)}{%
    \FormulaRefAuto{RingNegativeSum}[thm]{1,3,5}}
  \proofstepwidestarbreakkeep[1,2,3,4,5]{(p\oplus(-q))\oplus(r\oplus(-s))=(p\oplus r)\oplus(-(q\oplus s))}{%
    \rIE{13,12}}
\end{tabproofwide}
\endgroup

\FormulaThmDeltaK[Multiplikation zweier Differenzen]{%
  \begin{aligned}[t]&\RingAlg{R}{\oplus}{\odot}{0_R}{1_R}\dsep p,q,r,s\in R\vdash{}\\&(p\oplus(-q))\odot(r\oplus(-s))={}\\&\quad((p\odot r)\oplus(q\odot s))\oplus(-((p\odot s)\oplus(q\odot r)))\end{aligned}
}{RingDifferenceMultiplication}{%
  \DeltaRow{Mengen}{R\dsep0_R\dsep1_R\dsep p\dsep q\dsep r\dsep s}
  \DeltaRow{Funktionen}{\oplus\dsep\odot}
}
\begingroup
\ProofWideReasonsNearFormula
\begin{tabproofsplitwide}
  \proofpartwide{Träger}
  \proofstepwidestarbreakkeep[1]{\RingAlg{R}{\oplus}{\odot}{0_R}{1_R}}{%
    \rA}
  \proofstepwidestarbreakkeep[2]{p\in R}{%
    \rA}
  \proofstepwidestarbreakkeep[3]{q\in R}{%
    \rA}
  \proofstepwidestarbreakkeep[4]{r\in R}{%
    \rA}
  \proofstepwidestarbreakkeep[5]{s\in R}{%
    \rA}
  \proofstepwidestarbreakkeep[1]{\AbelianGroupAlg{R}{\oplus}{0_R}}{%
    \FormulaRefAuto{RingAdditiveAbelianGroupAxiom}[ax]{1}}
  \proofstepwidestarbreakkeep[1]{\GroupAlg{R}{\oplus}{0_R}}{%
    \FormulaRefAuto{RingAdditiveGroup}[thm]{1}}
  \proofstepwidestarbreakkeep[1]{\CommutativeMonoidAlg{R}{\oplus}{0_R}}{%
    \FormulaRefAuto{AbelianGroupIsCommutativeMonoid}[thm]{6}}
  \proofstepwidestarbreakkeep[1]{\CommutativeSemigroupAlg{R}{\oplus}}{%
    \FormulaRefAuto{CommutativeMonoidIsCommutativeSemigroup}[thm]{8}}
  \proofstepwidestarbreakkeep[1,3]{-q\in R}{%
    \FormulaRefAuto{RingNegativeCarrier}[thm]{1,3}}
  \proofstepwidestarbreakkeep[1,5]{-s\in R}{%
    \FormulaRefAuto{RingNegativeCarrier}[thm]{1,5}}
  \proofstepwidestarbreakkeep[1,2,3]{p\oplus(-q)\in R}{%
    \FormulaRefAuto{GroupClosure}[thm]{7,2,10}}
  \proofstepwidestarbreakkeep[1,2,4]{p\odot r\in R}{%
    \FormulaRefAuto{RingMultiplicationClosure}[thm]{1,2,4}}
  \proofstepwidestarbreakkeep[1,2,5]{p\odot s\in R}{%
    \FormulaRefAuto{RingMultiplicationClosure}[thm]{1,2,5}}
  \proofstepwidestarbreakkeep[1,3,4]{q\odot r\in R}{%
    \FormulaRefAuto{RingMultiplicationClosure}[thm]{1,3,4}}
  \proofstepwidestarbreakkeep[1,3,5]{q\odot s\in R}{%
    \FormulaRefAuto{RingMultiplicationClosure}[thm]{1,3,5}}
  \proofstepwidestarbreakkeep[1,2,5]{-(p\odot s)\in R}{%
    \FormulaRefAuto{RingNegativeCarrier}[thm]{1,14}}
  \proofstepwidestarbreakkeep[1,3,4]{-(q\odot r)\in R}{%
    \FormulaRefAuto{RingNegativeCarrier}[thm]{1,15}}
  \closeproofpartwide

  \proofpartwide{Distributivität und Vorzeichen}
  \setcounter{proofstepnr}{18}
  \proofstepwidestarbreakkeep[1,2,3,4,5]{(p\oplus(-q))\odot(r\oplus(-s))=((p\oplus(-q))\odot r)\oplus((p\oplus(-q))\odot(-s))}{%
    \FormulaRefAuto{RingLeftDistributivityAxiom}[ax]{1,12,4,11}}
  \proofstepwidestarbreakkeep[1,2,3,4]{(p\oplus(-q))\odot r=(p\odot r)\oplus((-q)\odot r)}{%
    \FormulaRefAuto{RingRightDistributivityAxiom}[ax]{1,2,10,4}}
  \proofstepwidestarbreakkeep[1,2,3,5]{(p\oplus(-q))\odot(-s)=(p\odot(-s))\oplus((-q)\odot(-s))}{%
    \FormulaRefAuto{RingRightDistributivityAxiom}[ax]{1,2,10,11}}
  \proofstepwidestarbreakkeep[1,3,4]{(-q)\odot r=-(q\odot r)}{%
    \FormulaRefAuto{RingLeftNegativeProduct}[thm]{1,3,4}}
  \proofstepwidestarbreakkeep[1,2,5]{p\odot(-s)=-(p\odot s)}{%
    \FormulaRefAuto{RingRightNegativeProduct}[thm]{1,2,5}}
  \proofstepwidestarbreakkeep[1,3,5]{(-q)\odot(-s)=q\odot s}{%
    \FormulaRefAuto{RingNegativeTimesNegative}[thm]{1,3,5}}
  \proofstepwidestarbreakkeep[1,2,3,4,5]{(p\oplus(-q))\odot(r\oplus(-s))=((p\odot r)\oplus(-(q\odot r)))\oplus((-(p\odot s))\oplus(q\odot s))}{%
    \rIE{20,21,22,23,24,19}}
  \closeproofpartwide

  \proofpartwide{Positive und negative Produkte zusammenfassen}
  \setcounter{proofstepnr}{25}
  \proofstepwidestarbreakkeep[1,2,3,5]{(-(p\odot s))\oplus(q\odot s)=(q\odot s)\oplus(-(p\odot s))}{%
    \FormulaRefAuto{AbelianGroupCommutativityAxiom}[ax]{6,17,16}}
  \proofstepwidestarbreakkeep[1,2,3,4,5]{((p\odot r)\oplus(-(q\odot r)))\oplus((q\odot s)\oplus(-(p\odot s)))=((p\odot r)\oplus(q\odot s))\oplus(-((q\odot r)\oplus(p\odot s)))}{%
    \FormulaRefAuto{RingDifferenceAddition}[thm]{1,13,15,16,14}}
  \proofstepwidestarbreakkeep[1,2,3,4,5]{(q\odot r)\oplus(p\odot s)=(p\odot s)\oplus(q\odot r)}{%
    \FormulaRefAuto{AbelianGroupCommutativityAxiom}[ax]{6,15,14}}
  \proofstepwidestarbreakkeep[1,2,3,4,5]{(p\oplus(-q))\odot(r\oplus(-s))=((p\odot r)\oplus(q\odot s))\oplus(-((p\odot s)\oplus(q\odot r)))}{%
    \rIE{26,28,27,25}}
  \closeproofpartwide
\end{tabproofsplitwide}
\endgroup

\chapter{Homomorphismen von Ringen}

\section{Der einem Ring zugrunde liegende Halbring}

Die Nullabsorption wurde für Ringe aus den übrigen Ringaxiomen hergeleitet.
Damit ist jeder Ring mit Eins insbesondere ein Halbring mit Eins. Dieses
Resultat erlaubt es, die kanonische Abbildung aus \(\mathbb N\) ohne eine
zweite Rekursionskonstruktion zu verwenden.

\FormulaThmDeltaK[Jeder Ring mit Eins ist ein Halbring mit Eins]{%
  \RingAlg{R}{+}{\cdot}{0}{1}
  \vdash\SemiringAlg{R}{+}{\cdot}{0}{1}%
}{RingUnderlyingSemiring}{%
  \DeltaRow{Mengen}{R\dsep0\dsep1\dsep a\dsep b\dsep c}
  \DeltaRow{Funktionen}{+\dsep\cdot}
}
\begin{tabproofsplitwide}
  \proofpartwide{Additive und multiplikative Monoidstruktur}
  \proofstepwidestar[1]{\RingAlg{R}{+}{\cdot}{0}{1}}{\rA}
  \proofstepwidestar[1]{\AbelianGroupAlg{R}{+}{0}}{%
    \FormulaRefAuto{RingAdditiveAbelianGroupAxiom}[ax]{1}}
  \proofstepwidestar[1]{\CommutativeMonoidAlg{R}{+}{0}}{%
    \FormulaRefAuto{AbelianGroupIsCommutativeMonoid}[thm]{2}}
  \proofstepwidestar[1]{\MonoidAlg{R}{\cdot}{1}}{%
    \FormulaRefAuto{RingMultiplicativeMonoidAxiom}[ax]{1}}
  \closeproofpartwide

  \proofpartwide{Distributivität}
  \setcounter{proofstepnr}{4}
  \proofstepwidestar[5]{a\in R}{\rA}
  \proofstepwidestar[6]{b\in R}{\rA}
  \proofstepwidestar[7]{c\in R}{\rA}
  \proofstepwidestar[1,5,6,7]{%
    a\cdot(b+c)=a\cdot b+a\cdot c}{%
    \FormulaRefAuto{RingLeftDistributivityAxiom}[ax]{1,5,6,7}}
  \proofstepwidestar[1,5,6,7]{%
    (a+b)\cdot c=a\cdot c+b\cdot c}{%
    \FormulaRefAuto{RingRightDistributivityAxiom}[ax]{1,5,6,7}}
  \closeproofpartwide

  \proofpartwide{Nullabsorption und Halbringstruktur}
  \setcounter{proofstepnr}{9}
  \proofstepwidestar[1,5]{0\cdot a=0}{%
    \FormulaRefAuto{RingLeftZeroAbsorption}[thm]{1,5}}
  \proofstepwidestar[1,5]{a\cdot0=0}{%
    \FormulaRefAuto{RingRightZeroAbsorption}[thm]{1,5}}
  \proofstepwidestar[1]{\SemiringAlg{R}{+}{\cdot}{0}{1}}{%
    \FormulaRefAuto{SemiringDef}[def]{3,4,8,9,10,11}}
  \closeproofpartwide
\end{tabproofsplitwide}

\section{Ringhomomorphismen}

Da alle Ringe dieses Bandes ein ausgezeichnetes Einselement besitzen, meint
\emph{Ringhomomorphismus} im Folgenden stets einen einserhaltenden
Ringhomomorphismus.

\FormulaDefDeltaK[Begriff des Ringhomomorphismus]{%
  \RingHom{f}{R,\oplus,\odot,0_R,1_R}
    {S,\boxplus,\boxtimes,0_S,1_S}%
}{RingHomDef}{%
  \DeltaRow{Mengen}{R\dsep S\dsep0_R\dsep1_R\dsep0_S\dsep1_S}
  \DeltaRow{Funktionen}{f\dsep\oplus\dsep\odot\dsep\boxplus\dsep\boxtimes}
  \DeltaRow{\textbf{Axiome}}{}
  \DeltaPrem{Quellring}{%
    \RingAlg{R}{\oplus}{\odot}{0_R}{1_R}}
  \DeltaPrem{Zielring}{%
    \RingAlg{S}{\boxplus}{\boxtimes}{0_S}{1_S}}
  \DeltaPrem{Additiver Anteil}{%
    \GroupHom{f}{R,\oplus,0_R}{S,\boxplus,0_S}}
  \DeltaPrem{Multiplikativer Anteil}{%
    \MonoidHom{f}{R,\odot,1_R}{S,\boxtimes,1_S}}
  \DeltaRow{\textbf{Neues Symbol}}{}
  \DeltaPrem{Symbol}{\mathsf{RingHom}}
}

\FormulaAxiomDeltaK[Trägerabbildung eines Ringhomomorphismus]{%
  \RingHom{f}{R,\oplus,\odot,0_R,1_R}
    {S,\boxplus,\boxtimes,0_S,1_S}
  \vdash f\colon R\to S%
}{RingHomFunctionAxiom}{%
  \DeltaRow{Mengen}{R\dsep S}
  \DeltaRow{Funktionen}{f\dsep\oplus\dsep\odot\dsep\boxplus\dsep\boxtimes}
}
\par\smallskip

\FormulaAxiomDeltaK[Additionserhaltung eines Ringhomomorphismus]{%
  \begin{aligned}[t]
    &\RingHom{f}{R,\oplus,\odot,0_R,1_R}
      {S,\boxplus,\boxtimes,0_S,1_S}\dsep x,y\in R\\
    &\vdash f(x\oplus y)=f(x)\boxplus f(y)
  \end{aligned}
}{RingHomAdditionAxiom}{%
  \DeltaRow{Mengen}{R\dsep S\dsep x\dsep y}
  \DeltaRow{Funktionen}{f\dsep\oplus\dsep\odot\dsep\boxplus\dsep\boxtimes}
}
\par\smallskip

\FormulaAxiomDeltaK[Multiplikationserhaltung eines Ringhomomorphismus]{%
  \begin{aligned}[t]
    &\RingHom{f}{R,\oplus,\odot,0_R,1_R}
      {S,\boxplus,\boxtimes,0_S,1_S}\dsep x,y\in R\\
    &\vdash f(x\odot y)=f(x)\boxtimes f(y)
  \end{aligned}
}{RingHomMultiplicationAxiom}{%
  \DeltaRow{Mengen}{R\dsep S\dsep x\dsep y}
  \DeltaRow{Funktionen}{f\dsep\oplus\dsep\odot\dsep\boxplus\dsep\boxtimes}
}
\par\smallskip

\FormulaAxiomDeltaK[Einserhaltung eines Ringhomomorphismus]{%
  \RingHom{f}{R,\oplus,\odot,0_R,1_R}
    {S,\boxplus,\boxtimes,0_S,1_S}
  \vdash f(1_R)=1_S%
}{RingHomOneAxiom}{%
  \DeltaRow{Mengen}{R\dsep S\dsep1_R\dsep1_S}
  \DeltaRow{Funktionen}{f\dsep\oplus\dsep\odot\dsep\boxplus\dsep\boxtimes}
}

\FormulaThmDeltaK[Ringhomomorphismen erhalten die Null]{%
  \RingHom{f}{R,\oplus,\odot,0_R,1_R}
    {S,\boxplus,\boxtimes,0_S,1_S}
  \vdash f(0_R)=0_S%
}{RingHomZeroPreservation}{%
  \DeltaRow{Mengen}{R\dsep S\dsep0_R\dsep0_S}
  \DeltaRow{Funktionen}{f\dsep\oplus\dsep\odot\dsep\boxplus\dsep\boxtimes}
}
\begin{tabproof}
  \proofstep{1}{%
    \RingHom{f}{R,\oplus,\odot,0_R,1_R}
      {S,\boxplus,\boxtimes,0_S,1_S}}{\rA}
  \proofstep{1}{\GroupHom{f}{R,\oplus,0_R}{S,\boxplus,0_S}}{%
    \FormulaRefAuto{RingHomDef}[def]{1}}
  \proofstep{1}{f(0_R)=0_S}{%
    \FormulaRefAuto{GroupHomIdentityPreservation}[thm]{2}}
\end{tabproof}

\FormulaThmDeltaK[Ringhomomorphismen erhalten additive Inverse]{%
  \RingHom{f}{R,\oplus,\odot,0_R,1_R}
    {S,\boxplus,\boxtimes,0_S,1_S}\dsep x\in R
  \vdash f(-x)=-f(x)%
}{RingHomNegativePreservation}{%
  \DeltaRow{Mengen}{R\dsep S\dsep x}
  \DeltaRow{Funktionen}{f\dsep\oplus\dsep\odot\dsep\boxplus\dsep\boxtimes}
}
\begin{tabproof}
  \proofstep{1}{%
    \RingHom{f}{R,\oplus,\odot,0_R,1_R}
      {S,\boxplus,\boxtimes,0_S,1_S}}{\rA}
  \proofstep{2}{x\in R}{\rA}
  \proofstep{1}{\GroupHom{f}{R,\oplus,0_R}{S,\boxplus,0_S}}{%
    \FormulaRefAuto{RingHomDef}[def]{1}}
  \proofstep{1,2}{f(-x)=-f(x)}{%
    \FormulaRefAuto{GroupHomInversePreservation}[thm]{3,2}}
\end{tabproof}

\section{Ringisomorphismen und Komposition}

\FormulaDefDeltaK[Begriff des Ringisomorphismus]{%
  \RingIso{f}{R,\oplus,\odot,0_R,1_R}
    {S,\boxplus,\boxtimes,0_S,1_S}%
}{RingIsoDef}{%
  \DeltaRow{Mengen}{R\dsep S}
  \DeltaRow{Funktionen}{f\dsep\oplus\dsep\odot\dsep\boxplus\dsep\boxtimes}
  \DeltaRow{\textbf{Axiome}}{}
  \DeltaPrem{Ringhomomorphismus}{%
    \RingHom{f}{R,\oplus,\odot,0_R,1_R}
      {S,\boxplus,\boxtimes,0_S,1_S}}
  \DeltaPrem{Bijektivität}{f\colon R\bij S}
  \DeltaRow{\textbf{Neues Symbol}}{}
  \DeltaPrem{Ringisomorphismus}{%
    \RingIso{f}{R,\oplus,\odot,0_R,1_R}
      {S,\boxplus,\boxtimes,0_S,1_S}}
}

\FormulaThmDeltaK[Komposition von Ringhomomorphismen]{%
  \begin{aligned}[t]
    &\RingHom{f}{R,\oplus,\odot,0_R,1_R}
      {S,\boxplus,\boxtimes,0_S,1_S}\dsep{}\\[-2pt]
    &\RingHom{g}{S,\boxplus,\boxtimes,0_S,1_S}
      {T,\uplus,\otimes,0_T,1_T}\\
    &\vdash
      \RingHom{g\circ f}{R,\oplus,\odot,0_R,1_R}
      {T,\uplus,\otimes,0_T,1_T}
  \end{aligned}
}{RingHomComposition}{%
  \DeltaRow{Mengen}{R\dsep S\dsep T}
  \DeltaRow{Funktionen}{f\dsep g\dsep\oplus\dsep\odot\dsep
    \boxplus\dsep\boxtimes\dsep\uplus\dsep\otimes}
}
\begin{tabproofsplitwide}
  \proofpartwide{Additiver Gruppenanteil}
  \proofstepwidestar[1]{%
    \RingHom{f}{R,\oplus,\odot,0_R,1_R}
      {S,\boxplus,\boxtimes,0_S,1_S}}{\rA}
  \proofstepwidestar[2]{%
    \RingHom{g}{S,\boxplus,\boxtimes,0_S,1_S}
      {T,\uplus,\otimes,0_T,1_T}}{\rA}
  \proofstepwidestar[1]{%
    \GroupHom{f}{R,\oplus,0_R}{S,\boxplus,0_S}}{%
    \FormulaRefAuto{RingHomDef}[def]{1}}
  \proofstepwidestar[2]{%
    \GroupHom{g}{S,\boxplus,0_S}{T,\uplus,0_T}}{%
    \FormulaRefAuto{RingHomDef}[def]{2}}
  \proofstepwidestar[1,2]{%
    \GroupHom{g\circ f}{R,\oplus,0_R}{T,\uplus,0_T}}{%
    \FormulaRefAuto{GroupHomComposition}[thm]{3,4}}
  \closeproofpartwide

  \proofpartwide{Multiplikativer Monoidanteil und Schluss}
  \setcounter{proofstepnr}{5}
  \proofstepwidestar[1]{%
    \MonoidHom{f}{R,\odot,1_R}{S,\boxtimes,1_S}}{%
    \FormulaRefAuto{RingHomDef}[def]{1}}
  \proofstepwidestar[2]{%
    \MonoidHom{g}{S,\boxtimes,1_S}{T,\otimes,1_T}}{%
    \FormulaRefAuto{RingHomDef}[def]{2}}
  \proofstepwidestar[1,2]{%
    \MonoidHom{g\circ f}{R,\odot,1_R}{T,\otimes,1_T}}{%
    \FormulaRefAuto{MonoidHomComposition}[thm]{6,7}}
  \proofstepwidestar[1]{%
    \RingAlg{R}{\oplus}{\odot}{0_R}{1_R}}{%
    \FormulaRefAuto{RingHomDef}[def]{1}}
  \proofstepwidestar[2]{%
    \RingAlg{T}{\uplus}{\otimes}{0_T}{1_T}}{%
    \FormulaRefAuto{RingHomDef}[def]{2}}
  \proofstepwidestar[1,2]{%
    \RingHom{g\circ f}{R,\oplus,\odot,0_R,1_R}
      {T,\uplus,\otimes,0_T,1_T}}{%
    \FormulaRefAuto{RingHomDef}[def]{9,10,5,8}}
  \closeproofpartwide
\end{tabproofsplitwide}

\chapter{Der Ring der ganzen Zahlen}

\section{Das multiplikative kommutative Monoid}

\FormulaThmDeltaK[Multiplikatives kommutatives Monoid der ganzen Zahlen]{%
  \CommutativeMonoidAlg{\mathbb Z}{\cdot}{\IntOne}%
}{IntegerMultiplicativeCommutativeMonoid}{%
  \DeltaRow{Mengen}{\mathbb Z\dsep \IntOne\dsep z\dsep w\dsep u}
  \DeltaRow{Funktionen}{\cdot}
}
\begin{tabproofsplitwide}
  \proofpartwideindR[Multiplikative Halbgruppe]{%
    \SemigroupAlg{\mathbb Z}{\cdot}%
  }{\SemigroupAlg{\mathbb Z}{\cdot}}[IntegerMultiplicativeSemigroup]
    \proofstepwidestar[1]{z\in\mathbb Z}{\rA}
    \proofstepwidestar[2]{w\in\mathbb Z}{\rA}
    \proofstepwidestar[1,2]{z\cdot w\in\mathbb Z}{%
      \FormulaRefAuto{IntegerMulClosureAxiom}[ax]{1,2}}
    \proofstepwidestar[4]{u\in\mathbb Z}{\rA}
    \proofstepwidestar[1,2,4]{%
      (z\cdot w)\cdot u=z\cdot(w\cdot u)}{%
      \FormulaRefAuto{IntegerMulAssociativeAxiom}[ax]{1,2,4}}
    \proofstepwidestar[]{\SemigroupAlg{\mathbb Z}{\cdot}}{%
      \FormulaRefAuto{SemigroupDef}[def]{3,5}}
  \closeproofpartwideindR

  \proofpartwideindR[Multiplikatives Monoid]{%
    \MonoidAlg{\mathbb Z}{\cdot}{\IntOne}%
  }{\MonoidAlg{\mathbb Z}{\cdot}{\IntOne}}[IntegerMultiplicativeMonoid]
    \proofstepwidestar[]{\SemigroupAlg{\mathbb Z}{\cdot}}{%
      \FormulaRefAuto{IntegerMultiplicativeSemigroup}[thm]}
    \proofstepwidestar[]{\IntOne\in\mathbb Z}{%
      \FormulaRefAuto{IntOneCarrier}[thm]}
    \proofstepwidestar[3]{z\in\mathbb Z}{\rA}
    \proofstepwidestar[3]{z\cdot\IntOne=z\land\IntOne\cdot z=z}{%
      \FormulaRefAuto{IntegerMulOneAxiom}[ax]{3}}
    \proofstepwidestar[3]{\IntOne\cdot z=z}{%
      \rAEb{4}}
    \proofstepwidestar[3]{z\cdot\IntOne=z\land\IntOne\cdot z=z}{%
      \FormulaRefAuto{IntegerMulOneAxiom}[ax]{3}}
    \proofstepwidestar[3]{z\cdot\IntOne=z}{%
      \rAEa{6}}
    \proofstepwidestar[]{\MonoidAlg{\mathbb Z}{\cdot}{\IntOne}}{%
      \FormulaRefAuto{MonoidDef}[def]{1,2,5,7}}
  \closeproofpartwideindR

  \proofpartwide{Kommutativität}
    \proofstepwidestar[]{\MonoidAlg{\mathbb Z}{\cdot}{\IntOne}}{%
      \FormulaRefAuto{IntegerMultiplicativeMonoid}[thm]}
    \proofstepwidestar[2]{z\in\mathbb Z}{\rA}
    \proofstepwidestar[3]{w\in\mathbb Z}{\rA}
    \proofstepwidestar[2,3]{z\cdot w=w\cdot z}{%
      \FormulaRefAuto{IntegerMulCommutativeAxiom}[ax]{2,3}}
    \proofstepwidestar[]{%
      \CommutativeMonoidAlg{\mathbb Z}{\cdot}{\IntOne}}{%
      \FormulaRefAuto{CommutativeMonoidDef}[def]{1,4}}
  \closeproofpartwide
\end{tabproofsplitwide}


\section{Die Ringstruktur}

\FormulaThmDeltaK[Ring mit Eins der ganzen Zahlen]{%
  \RingAlg{\mathbb Z}{+}{\cdot}{\IntZero}{\IntOne}
}{IntegerUnitalRing}{%
  \DeltaRow{Mengen}{\mathbb Z\dsep\IntZero\dsep\IntOne\dsep z\dsep w\dsep u}
  \DeltaRow{Funktionen}{+\dsep\cdot}
}
\begin{tabproofwide}
    \proofstepwidestar[]{%
      \AbelianGroupAlg{\mathbb Z}{+}{\IntZero}}{%
      \FormulaRefAuto{IntegerAdditiveAbelianGroup}[thm]}
    \proofstepwidestar[]{%
      \CommutativeMonoidAlg{\mathbb Z}{\cdot}{\IntOne}}{%
      \FormulaRefAuto{IntegerMultiplicativeCommutativeMonoid}[thm]}
    \proofstepwidestar[]{%
      \MonoidAlg{\mathbb Z}{\cdot}{\IntOne}}{%
      \FormulaRefAuto{CommutativeMonoidBaseMonoidAxiom}[ax]{2}}
    \proofstepwidestar[4]{z\in\mathbb Z}{\rA}
    \proofstepwidestar[5]{w\in\mathbb Z}{\rA}
    \proofstepwidestar[6]{u\in\mathbb Z}{\rA}
    \proofstepwidestar[4,5,6]{%
      z\cdot(w+u)=z\cdot w+z\cdot u}{%
      \FormulaRefAuto{IntegerMulLeftDistributiveAxiom}[ax]{4,5,6}}
    \proofstepwidestar[4,5,6]{%
      (z+w)\cdot u=z\cdot u+w\cdot u}{%
      \FormulaRefAuto{IntegerMulRightDistributiveAxiom}[ax]{4,5,6}}
    \proofstepwidestar[]{%
      \RingAlg{\mathbb Z}{+}{\cdot}{\IntZero}{\IntOne}}{%
      \FormulaRefAuto{RingDef}[def]{1,3,7,8}}
\end{tabproofwide}

\chapter{Der kanonische Ringhomomorphismus aus den ganzen Zahlen}

Sei \(\RingAlg{R}{\oplus}{\odot}{0_R}{1_R}\). Nach
\FormulaRefAuto{RingUnderlyingSemiring}[thm] ist \(R\) auch ein Halbring,
also steht die kanonische Abbildung
\(\NaturalSemiringMap{R,\oplus,\odot,0_R,1_R}\colon\mathbb N\to R\)
aus Band~39 zur Verfügung. Ein Differenzenpaar \((a,b)\) soll auf die
Differenz seiner beiden Bilder abgebildet werden.

\subsection{Repräsentantenunabhängigkeit}

\FormulaThmDeltaK[Kreuzsummen liefern gleiche Differenzen im Ring]{%
  \begin{aligned}[t]
    &\RingAlg{R}{\oplus}{\odot}{0_R}{1_R}\dsep
      p,q,r,s\in R\dsep p\oplus s=q\oplus r\\
    &\vdash p\oplus(-q)=r\oplus(-s)
  \end{aligned}
}{RingCrossSumGivesDifferenceEquality}{%
  \DeltaRow{Mengen}{R\dsep0_R\dsep1_R\dsep p\dsep q\dsep r\dsep s}
  \DeltaRow{Funktionen}{\oplus\dsep\odot}
}
\begingroup
\ProofWideReasonsNearFormula
\begin{tabproofsplitwide}
  \proofpartwide{Träger und additive Struktur}
  \proofstepwidestarbreakkeep[1]{\RingAlg{R}{\oplus}{\odot}{0_R}{1_R}}{%
    \rA}
  \proofstepwidestarbreakkeep[2]{p\in R}{%
    \rA}
  \proofstepwidestarbreakkeep[3]{q\in R}{%
    \rA}
  \proofstepwidestarbreakkeep[4]{r\in R}{%
    \rA}
  \proofstepwidestarbreakkeep[5]{s\in R}{%
    \rA}
  \proofstepwidestarbreakkeep[6]{p\oplus s=q\oplus r}{%
    \rA}
  \proofstepwidestarbreakkeep[1]{\AbelianGroupAlg{R}{\oplus}{0_R}}{%
    \FormulaRefAuto{RingAdditiveAbelianGroupAxiom}[ax]{1}}
  \proofstepwidestarbreakkeep[1]{\GroupAlg{R}{\oplus}{0_R}}{%
    \FormulaRefAuto{RingAdditiveGroup}[thm]{1}}
  \proofstepwidestarbreakkeep[1]{\CommutativeMonoidAlg{R}{\oplus}{0_R}}{%
    \FormulaRefAuto{AbelianGroupIsCommutativeMonoid}[thm]{7}}
  \proofstepwidestarbreakkeep[1]{\CommutativeSemigroupAlg{R}{\oplus}}{%
    \FormulaRefAuto{CommutativeMonoidIsCommutativeSemigroup}[thm]{9}}
  \proofstepwidestarbreakkeep[1,3]{-q\in R}{%
    \FormulaRefAuto{RingNegativeCarrier}[thm]{1,3}}
  \proofstepwidestarbreakkeep[1,5]{-s\in R}{%
    \FormulaRefAuto{RingNegativeCarrier}[thm]{1,5}}
  \proofstepwidestarbreakkeep[1,2,3]{p\oplus(-q)\in R}{%
    \FormulaRefAuto{GroupClosure}[thm]{8,2,11}}
  \proofstepwidestarbreakkeep[1,2,5]{p\oplus s\in R}{%
    \FormulaRefAuto{GroupClosure}[thm]{8,2,5}}
  \closeproofpartwide

  \proofpartwide{Einfügen und Umordnen der inversen Summanden}
  \setcounter{proofstepnr}{14}
  \proofstepwidestarbreakkeep[1,2,3]{(p\oplus(-q))\oplus0_R=p\oplus(-q)}{%
    \FormulaRefAuto{GroupRightIdentity}[thm]{8,13}}
  \proofstepwidestarbreakkeep[1,2,3]{p\oplus(-q)=(p\oplus(-q))\oplus0_R}{%
    \FormulaRefAuto{a=b\vdash b=a}[thm]{15}}
  \proofstepwidestarbreakkeep[1,5]{s\oplus(-s)=0_R}{%
    \FormulaRefAuto{RingNegativeRight}[thm]{1,5}}
  \proofstepwidestarbreakkeep[1,5]{(p\oplus(-q))\oplus0_R=(p\oplus(-q))\oplus(s\oplus(-s))}{%
    \rIE{17}}
  \proofstepwidestarbreakkeep[1,2,3,5]{(p\oplus(-q))\oplus(s\oplus(-s))=(p\oplus s)\oplus((-q)\oplus(-s))}{%
    \FormulaRefAuto{CommutativeSemigroupFourTermExchange}[thm]{10,2,11,5,12}}
  \proofstepwidestarbreakkeep[1,2,3,5]{((p\oplus s)\oplus(-q))\oplus(-s)=(p\oplus s)\oplus((-q)\oplus(-s))}{%
    \FormulaRefAuto{GroupAssociativity}[thm]{8,14,11,12}}
  \proofstepwidestarbreakkeep[1,2,3,5]{(p\oplus s)\oplus((-q)\oplus(-s))=((p\oplus s)\oplus(-q))\oplus(-s)}{%
    \FormulaRefAuto{a=b\vdash b=a}[thm]{20}}
  \closeproofpartwide

  \proofpartwide{Kreuzsumme einsetzen und kürzen}
  \setcounter{proofstepnr}{21}
  \proofstepwidestarbreakkeep[6]{((p\oplus s)\oplus(-q))\oplus(-s)=((q\oplus r)\oplus(-q))\oplus(-s)}{%
    \rIE{6}}
  \proofstepwidestarbreakkeep[1,3,4]{(q\oplus r)\oplus(-q)=(q\oplus(-q))\oplus r}{%
    \FormulaRefAuto{CommutativeSemigroupThreeTermExchange}[thm]{10,3,4,11}}
  \proofstepwidestarbreakkeep[1,3,4]{((q\oplus r)\oplus(-q))\oplus(-s)=((q\oplus(-q))\oplus r)\oplus(-s)}{%
    \rIE{23}}
  \proofstepwidestarbreakkeep[1,3]{q\oplus(-q)=0_R}{%
    \FormulaRefAuto{RingNegativeRight}[thm]{1,3}}
  \proofstepwidestarbreakkeep[1,3]{((q\oplus(-q))\oplus r)\oplus(-s)=(0_R\oplus r)\oplus(-s)}{%
    \rIE{25}}
  \proofstepwidestarbreakkeep[1,4]{0_R\oplus r=r}{%
    \FormulaRefAuto{GroupLeftIdentity}[thm]{8,4}}
  \proofstepwidestarbreakkeep[1,4]{(0_R\oplus r)\oplus(-s)=r\oplus(-s)}{%
    \rIE{27}}
  \proofstepwidestarbreakkeep[1,2,3,4,5,6]{p\oplus(-q)=r\oplus(-s)}{%
    \rChain{16,18,19,21,22,24,26,28}}
  \closeproofpartwide
\end{tabproofsplitwide}
\endgroup


\FormulaDefDeltaK[Wert eines Differenzenpaares im Zielring]{%
  \begin{aligned}[t]
    &\RingAlg{R}{\oplus}{\odot}{0_R}{1_R}\dsep a,b\in\mathbb N
    \vdash{}\\[-2pt]
    &\IntegerRingPairValue{R,\oplus,\odot,0_R,1_R}{a}{b}
      \coloneqq
      \NaturalSemiringMap{R,\oplus,\odot,0_R,1_R}(a)
      \oplus
      \bigl(
        -\NaturalSemiringMap{R,\oplus,\odot,0_R,1_R}(b)
      \bigr)
  \end{aligned}
}{IntegerRingPairValueDef}{%
  \DeltaRow{Mengen}{R\dsep0_R\dsep1_R\dsep a\dsep b}
  \DeltaRow{Funktionen}{\oplus\dsep\odot}
  \DeltaRow{\textbf{Neues Symbol}}{}
  \DeltaPrem{Wert eines Differenzenpaares}{%
    \IntegerRingPairValue{R,\oplus,\odot,0_R,1_R}{a}{b}}
}

\FormulaThmDeltaK[Der Wert eines Differenzenpaares liegt im Zielring]{%
  \RingAlg{R}{\oplus}{\odot}{0_R}{1_R}\dsep a,b\in\mathbb N
  \vdash
  \IntegerRingPairValue{R,\oplus,\odot,0_R,1_R}{a}{b}\in R%
}{IntegerRingPairValueCarrier}{%
  \DeltaRow{Mengen}{R\dsep a\dsep b}
  \DeltaRow{Funktionen}{\oplus\dsep\odot}
}
Für diesen Beweis setzen wir
\[
  \begin{aligned}
    \eta_R(n)
      &:=\NaturalSemiringMap{R,\oplus,\odot,0_R,1_R}(n),\\
    v_R(a,b)
      &:=\IntegerRingPairValue{R,\oplus,\odot,0_R,1_R}{a}{b}.
  \end{aligned}
\]
\begin{tabproof}
  \proofstep{1}{\RingAlg{R}{\oplus}{\odot}{0_R}{1_R}}{\rA}
  \proofstep{2}{a\in\mathbb N}{\rA}
  \proofstep{3}{b\in\mathbb N}{\rA}
  \proofstep{1}{\SemiringAlg{R}{\oplus}{\odot}{0_R}{1_R}}{%
    \FormulaRefAuto{RingUnderlyingSemiring}[thm]{1}}
  \proofstep{1}{\eta_R\colon\mathbb N\to R}{%
    \FormulaRefAuto{NaturalSemiringMapFunction}[thm]{4}}
  \proofstep{1,2}{\eta_R(a)\in R}{%
    \FormulaRefAuto{F\colon A\to B\dsep a\in A\vdash F(a)\in B}{5,2}}
  \proofstep{1,3}{\eta_R(b)\in R}{%
    \FormulaRefAuto{F\colon A\to B\dsep a\in A\vdash F(a)\in B}{5,3}}
  \proofstep{1,3}{-\eta_R(b)\in R}{%
    \FormulaRefAuto{RingNegativeCarrier}[thm]{1,7}}
  \proofstep{1,2,3}{%
    \eta_R(a)\oplus\bigl(-\eta_R(b)\bigr)\in R}{%
    \FormulaRefAuto{GroupClosure}[thm]{%
      \FormulaRefAuto{RingAdditiveGroup}[thm]{1},6,8}}
  \proofstep{1,2,3}{v_R(a,b)=\eta_R(a)\oplus\bigl(-\eta_R(b)\bigr)}{%
    \FormulaRefAuto{IntegerRingPairValueDef}[def]{1,2,3}}
  \proofstep{1,2,3}{v_R(a,b)\in R}{%
    \rIE{10,9}}
\end{tabproof}

\FormulaThmDeltaK[Repräsentantenunabhängigkeit des Ganzzahlbildes]{%
  \begin{aligned}[t]
    &\RingAlg{R}{\oplus}{\odot}{0_R}{1_R}\dsep
      a,b,c,d\in\mathbb N\dsep
      \IntClass{a}{b}=\IntClass{c}{d}\\
    &\vdash
      \IntegerRingPairValue{R,\oplus,\odot,0_R,1_R}{a}{b}
      =
      \IntegerRingPairValue{R,\oplus,\odot,0_R,1_R}{c}{d}
  \end{aligned}
}{IntegerRingMapRepresentativeIndependent}{%
  \DeltaRow{Mengen}{R\dsep a\dsep b\dsep c\dsep d}
  \DeltaRow{Funktionen}{\oplus\dsep\odot}
}
Zur Entlastung der Beweiszeilen schreiben wir hier
\[
  \begin{aligned}
    \eta_R(n)
      &:=\NaturalSemiringMap{R,\oplus,\odot,0_R,1_R}(n),\\
    v_R(x,y)
      &:=\IntegerRingPairValue{R,\oplus,\odot,0_R,1_R}{x}{y}.
  \end{aligned}
\]
\begingroup
\ProofWideReasonsNearFormula
\begin{tabproofsplitwide}
  \proofpartwide{Träger der natürlichen Bilder}
  \proofstepwidestarbreakkeep[1]{\RingAlg{R}{\oplus}{\odot}{0_R}{1_R}}{%
    \rA}
  \proofstepwidestarbreakkeep[2]{a\in\mathbb N}{%
    \rA}
  \proofstepwidestarbreakkeep[3]{b\in\mathbb N}{%
    \rA}
  \proofstepwidestarbreakkeep[4]{c\in\mathbb N}{%
    \rA}
  \proofstepwidestarbreakkeep[5]{d\in\mathbb N}{%
    \rA}
  \proofstepwidestarbreakkeep[1]{\SemiringAlg{R}{\oplus}{\odot}{0_R}{1_R}}{%
    \FormulaRefAuto{RingUnderlyingSemiring}[thm]{1}}
  \proofstepwidestarbreakkeep[1]{\eta_R\colon\mathbb N\to R}{%
    \FormulaRefAuto{NaturalSemiringMapFunction}[thm]{6}}
  \proofstepwidestarbreakkeep[1,2]{\eta_R(a)\in R}{%
    \FormulaRefAuto{F\colon A\to B\dsep a\in A\vdash F(a)\in B}[thm]{7,2}}
  \proofstepwidestarbreakkeep[1,3]{\eta_R(b)\in R}{%
    \FormulaRefAuto{F\colon A\to B\dsep a\in A\vdash F(a)\in B}[thm]{7,3}}
  \proofstepwidestarbreakkeep[1,4]{\eta_R(c)\in R}{%
    \FormulaRefAuto{F\colon A\to B\dsep a\in A\vdash F(a)\in B}[thm]{7,4}}
  \proofstepwidestarbreakkeep[1,5]{\eta_R(d)\in R}{%
    \FormulaRefAuto{F\colon A\to B\dsep a\in A\vdash F(a)\in B}[thm]{7,5}}
  \proofstepwidestarbreakkeep[12]{\IntClass{a}{b}=\IntClass{c}{d}}{%
    \rA}
  \proofstepwidestarbreakkeep[2,3,4,5,12]{a+d=b+c}{%
    \FormulaRefAuto{IntegerClassEquality}[thm]{2,3,4,5,12}}
  \closeproofpartwide

  \proofpartwide{Kreuzsummen im Zielring}
  \setcounter{proofstepnr}{13}
  \proofstepwidestarbreakkeep[1,2,5]{\eta_R(a+d)=\eta_R(a)\oplus\eta_R(d)}{%
    \FormulaRefAuto{NaturalSemiringMapAddition}[thm]{6,2,5}}
  \proofstepwidestarbreakkeep[1,3,4]{\eta_R(b+c)=\eta_R(b)\oplus\eta_R(c)}{%
    \FormulaRefAuto{NaturalSemiringMapAddition}[thm]{6,3,4}}
  \proofstepwidestarbreakkeep[2,3,4,5,12]{\eta_R(a+d)=\eta_R(b+c)}{%
    \rIE{13}}
  \proofstepwidestarbreakkeep[1,2,3,4,5,12]{\eta_R(a)\oplus\eta_R(d)=\eta_R(b)\oplus\eta_R(c)}{%
    \rIE{14,15,16}}
  \proofstepwidestarbreakkeep[1,2,3,4,5,12]{\eta_R(a)\oplus(-\eta_R(b))=\eta_R(c)\oplus(-\eta_R(d))}{%
    \FormulaRefAuto{RingCrossSumGivesDifferenceEquality}[thm]{1,8,9,10,11,17}}
  \proofstepwidestarbreakkeep[1,2,3]{v_R(a,b)=\eta_R(a)\oplus(-\eta_R(b))}{%
    \FormulaRefAuto{IntegerRingPairValueDef}[def]{1,2,3}}
  \proofstepwidestarbreakkeep[1,4,5]{v_R(c,d)=\eta_R(c)\oplus(-\eta_R(d))}{%
    \FormulaRefAuto{IntegerRingPairValueDef}[def]{1,4,5}}
  \proofstepwidestarbreakkeep[1,2,3,4,5,12]{v_R(a,b)=v_R(c,d)}{%
    \rIE{19,20,18}}
  \closeproofpartwide
\end{tabproofsplitwide}
\endgroup


\subsection{Quotientenabstieg}

\FormulaThmDeltaK[Eindeutiger Quotientenabstieg des Ganzzahlbildes]{%
  \begin{aligned}[t]
    &\RingAlg{R}{\oplus}{\odot}{0_R}{1_R}\vdash{}\\[-2pt]
    &\exists! f\Bigl(
      f\colon\mathbb Z\to R\land
      \forall a,b\in\mathbb N\,
      f(\IntClass{a}{b})
      =
      \IntegerRingPairValue{R,\oplus,\odot,0_R,1_R}{a}{b}
    \Bigr)
  \end{aligned}
}{IntegerRingMapQuotientDescent}{%
  \DeltaRow{Mengen}{%
    R\dsep a\dsep b\dsep c\dsep d\dsep z\dsep u\dsep u_1\dsep u_2}
  \DeltaRow{Funktionen}{f\dsep g\dsep\oplus\dsep\odot}
  \DeltaRow{Repräsentanten}{%
    \FormulaRefAuto{IntegerPairRepresentative}[thm]}
  \DeltaRow{Werte im Zielring}{%
    \FormulaRefAuto{IntegerRingPairValueCarrier}[thm]}
  \DeltaRow{Repräsentantenunabhängigkeit}{%
    \FormulaRefAuto{IntegerRingMapRepresentativeIndependent}[thm]}
}

Für den folgenden Beweis setzen wir lokal
\[
  \begin{aligned}
    v_R(a,b)&:=\IntegerRingPairValue{R,\oplus,\odot,0_R,1_R}{a}{b},\\
    \kappa(a,b)&:=\IntClass{a}{b},\\
    P(a,b)&:\Longleftrightarrow a\in\mathbb N\land b\in\mathbb N,\\
    E(f)&:\Longleftrightarrow
      \forall a\in\mathbb N\,\forall b\in\mathbb N\,
        f(\kappa(a,b))=v_R(a,b),\\
    D(f)&:\Longleftrightarrow
      \forall a\forall b\,
        \bigl(P(a,b)\rightarrow f(\kappa(a,b))=v_R(a,b)\bigr).
  \end{aligned}
\]
Der allgemeine Abstieg konstruiert die Abbildung aus diesen
Repräsentantenwerten. Die anschließende Quantorenumformung führt seine
Aussage in die Schreibweise des Satzes zurück.
\begin{tabproofsplitwide}
  \proofpartwide{Repräsentation der ganzen Zahlen und Träger der Werte}
    \proofstepwidestar[1]{%
      \RingAlg{R}{\oplus}{\odot}{0_R}{1_R}}{\rA}
    \proofstepwidestar[2]{P(a,b)}{\rA}
    \proofstepwidestar[2]{\kappa(a,b)\in\mathbb Z}{%
      \FormulaRefAuto{IntegerClassInZ}[thm]{\rAEa{2},\rAEb{2}}}
    \proofstepwidestar[]{%
      \forall a\forall b\,(P(a,b)\rightarrow\kappa(a,b)\in\mathbb Z)}{%
      \rUI{\rUI{\rRI{2,3}}}}
    \proofstepwidestar[5]{z\in\mathbb Z}{\rA}
    \proofstepwidestar[5]{%
      \exists a\in\mathbb N\,\exists b\in\mathbb N\,z=\kappa(a,b)}{%
      \FormulaRefAuto{IntegerPairRepresentative}[thm]{5}}
    \proofstepwidestar[]{%
      \begin{aligned}[t]
        &\bigl(\exists a\in\mathbb N\,\exists b\in\mathbb N\,
          z=\kappa(a,b)\bigr)\\[-2pt]
        &\quad\leftrightarrow\exists a\exists b\,
          (P(a,b)\land z=\kappa(a,b))
      \end{aligned}}{\FormulaRefAuto{BoundedPairExistsGuard}[thm]}
    \proofstepwidestar[5]{%
      \exists a\exists b\,(P(a,b)\land z=\kappa(a,b))}{%
      \rRE{\rLREa{7},6}}
    \proofstepwidestar[]{%
      \forall z\in\mathbb Z\,\exists a\exists b\,
        (P(a,b)\land z=\kappa(a,b))}{%
      \rUI{\rRI{5,8}}}
    \proofstepwidestar[10]{P(a,b)}{\rA}
    \proofstepwidestar[1,10]{v_R(a,b)\in R}{%
      \FormulaRefAuto{IntegerRingPairValueCarrier}[thm]{1,\rAEa{10},\rAEb{10}}}
    \proofstepwidestar[1]{%
      \forall a\forall b\,(P(a,b)\rightarrow v_R(a,b)\in R)}{%
      \rUI{\rUI{\rRI{10,11}}}}
  \closeproofpartwide

  \proofpartwide{Faserverträglichkeit und eindeutiger Abstieg}
    \setcounter{proofstepnr}{12}
    \proofstepwidestar[13]{%
      P(a,b)\land P(c,d)\land\kappa(a,b)=\kappa(c,d)}{\rA}
    \proofstepwidestar[13]{P(a,b)}{\rAEa{13}}
    \proofstepwidestar[13]{P(c,d)}{\rAEa{\rAEb{13}}}
    \proofstepwidestar[13]{\kappa(a,b)=\kappa(c,d)}{\rAEb{\rAEb{13}}}
    \proofstepwidestar[1,13]{v_R(a,b)=v_R(c,d)}{%
      \FormulaRefAuto{IntegerRingMapRepresentativeIndependent}[thm]{%
        1,\rAEa{14},\rAEb{14},\rAEa{15},\rAEb{15},16}}
    \proofstepwidestar[1]{%
      \begin{aligned}[t]
        &\forall a\forall b\forall c\forall d\,\Bigl(\\[-2pt]
        &\quad\bigl(P(a,b)\land P(c,d)\land\kappa(a,b)=\kappa(c,d)\bigr)\\[-2pt]
        &\quad\rightarrow v_R(a,b)=v_R(c,d)\Bigr)
      \end{aligned}}{%
      \rUI{\rUI{\rUI{\rUI{\rRI{13,17}}}}}}
    \proofstepwidestar[1]{%
      \exists! f\,\bigl(f\colon\mathbb Z\to R\land D(f)\bigr)}{%
      \FormulaRefAuto{RepresentativePairFactorization}[thm]{4,9,12,18}}
  \closeproofpartwide

  \proofpartwide{Rückkehr zu den beschränkten Quantoren}
    \setcounter{proofstepnr}{19}
    \proofstepwidestar[]{E(f)\leftrightarrow D(f)}{%
      \FormulaRefAuto{BoundedPairForallGuard}[thm]}
    \proofstepwidestar[]{D(f)\leftrightarrow E(f)}{%
      \FormulaRefAuto{P \leftrightarrow Q \vdash Q \leftrightarrow P}[thm]{20}}
    \proofstepwidestar[]{\forall f\,(D(f)\leftrightarrow E(f))}{\rUI{21}}
    \proofstepwidestar[1]{\exists! f\,\bigl(f\colon\mathbb Z\to R\land E(f)\bigr)}{%
      \FormulaRefAuto{UniqueExistenceUnderEquivalentConjuncts}[thm]{22,19}}
  \closeproofpartwide
\end{tabproofsplitwide}

\FormulaDefDeltaK[Kanonische Abbildung aus den ganzen Zahlen]{%
  \begin{aligned}[t]
    &\RingAlg{R}{\oplus}{\odot}{0_R}{1_R}\vdash{}\\[-2pt]
    &\IntegerRingMap{R,\oplus,\odot,0_R,1_R}
      \colon\mathbb Z\to R,\\[-2pt]
    &\forall a,b\in\mathbb N\,
      \Bigl(
        \IntegerRingMap{R,\oplus,\odot,0_R,1_R}(\IntClass{a}{b})
        \coloneqq
        \IntegerRingPairValue{R,\oplus,\odot,0_R,1_R}{a}{b}
      \Bigr)
  \end{aligned}
}{IntegerRingMapDef}{%
  \DeltaRow{Mengen}{R\dsep a\dsep b}
  \DeltaRow{Funktionen}{\oplus\dsep\odot}
  \DeltaRow{Repräsentanten}{%
    \FormulaRefAuto{IntegerPairRepresentative}[thm]}
  \DeltaRow{Repräsentantenunabhängigkeit}{%
    \FormulaRefAuto{IntegerRingMapRepresentativeIndependent}[thm]}
  \DeltaRow{Existenz und Eindeutigkeit}{%
    \FormulaRefAuto{IntegerRingMapQuotientDescent}[thm]}
  \DeltaRow{\textbf{Neues Symbol}}{}
  \DeltaPrem{Kanonische Abbildung}{%
    \IntegerRingMap{R,\oplus,\odot,0_R,1_R}}
}

\FormulaThmDeltaK[Trägerabbildung der kanonischen Ganzzahlabbildung]{%
  \RingAlg{R}{\oplus}{\odot}{0_R}{1_R}
  \vdash
  \IntegerRingMap{R,\oplus,\odot,0_R,1_R}
  \colon\mathbb Z\to R%
}{IntegerRingMapFunction}{%
  \DeltaRow{Mengen}{R}
  \DeltaRow{Funktionen}{\oplus\dsep\odot}
}
\begin{tabproof}
  \proofstep{1}{\RingAlg{R}{\oplus}{\odot}{0_R}{1_R}}{\rA}
  \proofstep{1}{%
    \IntegerRingMap{R,\oplus,\odot,0_R,1_R}
    \colon\mathbb Z\to R}{%
    \FormulaRefAuto{IntegerRingMapDef}[def]{1}}
\end{tabproof}
\par\smallskip

\FormulaThmDeltaK[Auswertung an einer Ganzzahlklasse]{%
  \begin{aligned}[t]
    &\RingAlg{R}{\oplus}{\odot}{0_R}{1_R}\dsep a,b\in\mathbb N\\
    &\vdash
      \IntegerRingMap{R,\oplus,\odot,0_R,1_R}(\IntClass{a}{b})
      =
      \NaturalSemiringMap{R,\oplus,\odot,0_R,1_R}(a)
      \oplus
      \bigl(-\NaturalSemiringMap{R,\oplus,\odot,0_R,1_R}(b)\bigr)
  \end{aligned}
}{IntegerRingMapClassValue}{%
  \DeltaRow{Mengen}{R\dsep a\dsep b}
  \DeltaRow{Funktionen}{\oplus\dsep\odot}
}
\[
  \eta_R=\NaturalSemiringMap{R,\oplus,\odot,0_R,1_R},\qquad
  \Phi_R=\IntegerRingMap{R,\oplus,\odot,0_R,1_R},
  \qquad
  v_R(a,b)=\IntegerRingPairValue{R,\oplus,\odot,0_R,1_R}{a}{b}.
\]
\begingroup
\ProofWideReasonsNearFormula
\begin{tabproofwide}
  \proofstepwidestarbreakkeep[1]{\RingAlg{R}{\oplus}{\odot}{0_R}{1_R}}{%
    \rA}
  \proofstepwidestarbreakkeep[2]{a\in\mathbb N}{%
    \rA}
  \proofstepwidestarbreakkeep[3]{b\in\mathbb N}{%
    \rA}
  \proofstepwidestarbreakkeep[1,2,3]{\Phi_R(\IntClass{a}{b})=v_R(a,b)}{%
    \FormulaRefAuto{IntegerRingMapDef}[def]{1,2,3}}
  \proofstepwidestarbreakkeep[1,2,3]{v_R(a,b)=\eta_R(a)\oplus(-\eta_R(b))}{%
    \FormulaRefAuto{IntegerRingPairValueDef}[def]{1,2,3}}
  \proofstepwidestarbreakkeep[1,2,3]{\Phi_R(\IntClass{a}{b})=\eta_R(a)\oplus(-\eta_R(b))}{%
    \rChain{4,5}}
\end{tabproofwide}
\endgroup


\subsection{Erhaltung der Ringoperationen}

\FormulaThmDeltaK[Addition der Differenzenpaarwerte]{%
  \begin{aligned}[t]&\RingAlg{R}{\oplus}{\odot}{0_R}{1_R}\dsep a,b,c,d\in\mathbb N\vdash{}\\&\IntegerRingPairValue{R,\oplus,\odot,0_R,1_R}{a+c}{b+d}={}\\&\quad\IntegerRingPairValue{R,\oplus,\odot,0_R,1_R}{a}{b}\oplus\IntegerRingPairValue{R,\oplus,\odot,0_R,1_R}{c}{d}\end{aligned}
}{IntegerRingPairValueAddition}{%
  \DeltaRow{Mengen}{R\dsep0_R\dsep1_R\dsep a\dsep b\dsep c\dsep d}
  \DeltaRow{Funktionen}{\oplus\dsep\odot}
}
\[
  \eta_R=\NaturalSemiringMap{R,\oplus,\odot,0_R,1_R},\qquad
  v_R(a,b)=\IntegerRingPairValue{R,\oplus,\odot,0_R,1_R}{a}{b}.
\]
\begingroup
\ProofWideReasonsNearFormula
\begin{tabproofsplitwide}
  \proofpartwide{Träger und natürliche Abbildung}
  \proofstepwidestarbreakkeep[1]{\RingAlg{R}{\oplus}{\odot}{0_R}{1_R}}{%
    \rA}
  \proofstepwidestarbreakkeep[2]{a\in\mathbb N}{%
    \rA}
  \proofstepwidestarbreakkeep[3]{b\in\mathbb N}{%
    \rA}
  \proofstepwidestarbreakkeep[4]{c\in\mathbb N}{%
    \rA}
  \proofstepwidestarbreakkeep[5]{d\in\mathbb N}{%
    \rA}
  \proofstepwidestarbreakkeep[1]{\SemiringAlg{R}{\oplus}{\odot}{0_R}{1_R}}{%
    \FormulaRefAuto{RingUnderlyingSemiring}[thm]{1}}
  \proofstepwidestarbreakkeep[1]{\eta_R\colon\mathbb N\to R}{%
    \FormulaRefAuto{NaturalSemiringMapFunction}[thm]{6}}
  \proofstepwidestarbreakkeep[1,2]{\eta_R(a)\in R}{%
    \FormulaRefAuto{F\colon A\to B\dsep a\in A\vdash F(a)\in B}[thm]{7,2}}
  \proofstepwidestarbreakkeep[1,3]{\eta_R(b)\in R}{%
    \FormulaRefAuto{F\colon A\to B\dsep a\in A\vdash F(a)\in B}[thm]{7,3}}
  \proofstepwidestarbreakkeep[1,4]{\eta_R(c)\in R}{%
    \FormulaRefAuto{F\colon A\to B\dsep a\in A\vdash F(a)\in B}[thm]{7,4}}
  \proofstepwidestarbreakkeep[1,5]{\eta_R(d)\in R}{%
    \FormulaRefAuto{F\colon A\to B\dsep a\in A\vdash F(a)\in B}[thm]{7,5}}
  \proofstepwidestarbreakkeep[2,4]{a+c\in\mathbb N}{%
    \FormulaRefAuto{PeanoAddClosure}[thm]{2,4}}
  \proofstepwidestarbreakkeep[3,5]{b+d\in\mathbb N}{%
    \FormulaRefAuto{PeanoAddClosure}[thm]{3,5}}
  \closeproofpartwide

  \proofpartwide{Paarwert der verrechneten Repräsentanten}
  \setcounter{proofstepnr}{13}
  \proofstepwidestarbreakkeep[1,2,4]{\eta_R(a+c)=\eta_R(a)\oplus\eta_R(c)}{%
    \FormulaRefAuto{NaturalSemiringMapAddition}[thm]{6,2,4}}
  \proofstepwidestarbreakkeep[1,3,5]{\eta_R(b+d)=\eta_R(b)\oplus\eta_R(d)}{%
    \FormulaRefAuto{NaturalSemiringMapAddition}[thm]{6,3,5}}
  \proofstepwidestarbreakkeep[1,2,3,4,5]{v_R(a+c,b+d)=\eta_R(a+c)\oplus(-\eta_R(b+d))}{%
    \FormulaRefAuto{IntegerRingPairValueDef}[def]{1,12,13}}
  \proofstepwidestarbreakkeep[1,2,3]{v_R(a,b)=\eta_R(a)\oplus(-\eta_R(b))}{%
    \FormulaRefAuto{IntegerRingPairValueDef}[def]{1,2,3}}
  \proofstepwidestarbreakkeep[1,4,5]{v_R(c,d)=\eta_R(c)\oplus(-\eta_R(d))}{%
    \FormulaRefAuto{IntegerRingPairValueDef}[def]{1,4,5}}
  \proofstepwidestarbreakkeep[1,2,3,4,5]{(\eta_R(a)\oplus(-\eta_R(b)))\oplus(\eta_R(c)\oplus(-\eta_R(d)))=(\eta_R(a)\oplus\eta_R(c))\oplus(-(\eta_R(b)\oplus\eta_R(d)))}{%
    \FormulaRefAuto{RingDifferenceAddition}[thm]{1,8,9,10,11}}
  \proofstepwidestarbreakkeep[1,2,3,4,5]{v_R(a+c,b+d)=v_R(a,b)\oplus v_R(c,d)}{%
    \rIE{14,15,17,18,19,16}}
  \closeproofpartwide
\end{tabproofsplitwide}
\endgroup

\FormulaThmDeltaK[Multiplikation der Differenzenpaarwerte]{%
  \begin{aligned}[t]&\RingAlg{R}{\oplus}{\odot}{0_R}{1_R}\dsep a,b,c,d\in\mathbb N\vdash{}\\&\IntegerRingPairValue{R,\oplus,\odot,0_R,1_R}{ac+bd}{ad+bc}={}\\&\quad\IntegerRingPairValue{R,\oplus,\odot,0_R,1_R}{a}{b}\odot\IntegerRingPairValue{R,\oplus,\odot,0_R,1_R}{c}{d}\end{aligned}
}{IntegerRingPairValueMultiplication}{%
  \DeltaRow{Mengen}{R\dsep0_R\dsep1_R\dsep a\dsep b\dsep c\dsep d}
  \DeltaRow{Funktionen}{\oplus\dsep\odot}
}
\[
  \eta_R=\NaturalSemiringMap{R,\oplus,\odot,0_R,1_R},\qquad
  v_R(a,b)=\IntegerRingPairValue{R,\oplus,\odot,0_R,1_R}{a}{b}.
\]
\begingroup
\ProofWideReasonsNearFormula
\begin{tabproofsplitwide}
  \proofpartwide{Träger und natürliche Abbildung}
  \proofstepwidestarbreakkeep[1]{\RingAlg{R}{\oplus}{\odot}{0_R}{1_R}}{%
    \rA}
  \proofstepwidestarbreakkeep[2]{a\in\mathbb N}{%
    \rA}
  \proofstepwidestarbreakkeep[3]{b\in\mathbb N}{%
    \rA}
  \proofstepwidestarbreakkeep[4]{c\in\mathbb N}{%
    \rA}
  \proofstepwidestarbreakkeep[5]{d\in\mathbb N}{%
    \rA}
  \proofstepwidestarbreakkeep[1]{\SemiringAlg{R}{\oplus}{\odot}{0_R}{1_R}}{%
    \FormulaRefAuto{RingUnderlyingSemiring}[thm]{1}}
  \proofstepwidestarbreakkeep[1]{\eta_R\colon\mathbb N\to R}{%
    \FormulaRefAuto{NaturalSemiringMapFunction}[thm]{6}}
  \proofstepwidestarbreakkeep[1,2]{\eta_R(a)\in R}{%
    \FormulaRefAuto{F\colon A\to B\dsep a\in A\vdash F(a)\in B}[thm]{7,2}}
  \proofstepwidestarbreakkeep[1,3]{\eta_R(b)\in R}{%
    \FormulaRefAuto{F\colon A\to B\dsep a\in A\vdash F(a)\in B}[thm]{7,3}}
  \proofstepwidestarbreakkeep[1,4]{\eta_R(c)\in R}{%
    \FormulaRefAuto{F\colon A\to B\dsep a\in A\vdash F(a)\in B}[thm]{7,4}}
  \proofstepwidestarbreakkeep[1,5]{\eta_R(d)\in R}{%
    \FormulaRefAuto{F\colon A\to B\dsep a\in A\vdash F(a)\in B}[thm]{7,5}}
  \proofstepwidestarbreakkeep[2,4]{ac\in\mathbb N}{%
    \FormulaRefAuto{PeanoMulClosure}[thm]{2,4}}
  \proofstepwidestarbreakkeep[1,2,4]{\eta_R(ac)=\eta_R(a)\odot\eta_R(c)}{%
    \FormulaRefAuto{NaturalSemiringMapMultiplication}[thm]{6,2,4}}
  \proofstepwidestarbreakkeep[3,5]{bd\in\mathbb N}{%
    \FormulaRefAuto{PeanoMulClosure}[thm]{3,5}}
  \proofstepwidestarbreakkeep[1,3,5]{\eta_R(bd)=\eta_R(b)\odot\eta_R(d)}{%
    \FormulaRefAuto{NaturalSemiringMapMultiplication}[thm]{6,3,5}}
  \proofstepwidestarbreakkeep[2,5]{ad\in\mathbb N}{%
    \FormulaRefAuto{PeanoMulClosure}[thm]{2,5}}
  \proofstepwidestarbreakkeep[1,2,5]{\eta_R(ad)=\eta_R(a)\odot\eta_R(d)}{%
    \FormulaRefAuto{NaturalSemiringMapMultiplication}[thm]{6,2,5}}
  \proofstepwidestarbreakkeep[3,4]{bc\in\mathbb N}{%
    \FormulaRefAuto{PeanoMulClosure}[thm]{3,4}}
  \proofstepwidestarbreakkeep[1,3,4]{\eta_R(bc)=\eta_R(b)\odot\eta_R(c)}{%
    \FormulaRefAuto{NaturalSemiringMapMultiplication}[thm]{6,3,4}}
  \proofstepwidestarbreakkeep[2,3,4,5]{ac+bd\in\mathbb N}{%
    \FormulaRefAuto{PeanoAddClosure}[thm]{12,14}}
  \proofstepwidestarbreakkeep[2,3,4,5]{ad+bc\in\mathbb N}{%
    \FormulaRefAuto{PeanoAddClosure}[thm]{16,18}}
  \closeproofpartwide

  \proofpartwide{Paarwert der verrechneten Repräsentanten}
  \setcounter{proofstepnr}{21}
  \proofstepwidestarbreakkeep[1,2,3,4,5]{\eta_R(ac+bd)=\eta_R(ac)\oplus\eta_R(bd)}{%
    \FormulaRefAuto{NaturalSemiringMapAddition}[thm]{6,12,14}}
  \proofstepwidestarbreakkeep[1,2,3,4,5]{\eta_R(ad+bc)=\eta_R(ad)\oplus\eta_R(bc)}{%
    \FormulaRefAuto{NaturalSemiringMapAddition}[thm]{6,16,18}}
  \proofstepwidestarbreakkeep[1,2,3,4,5]{v_R(ac+bd,ad+bc)=\eta_R(ac+bd)\oplus(-\eta_R(ad+bc))}{%
    \FormulaRefAuto{IntegerRingPairValueDef}[def]{1,20,21}}
  \proofstepwidestarbreakkeep[1,2,3]{v_R(a,b)=\eta_R(a)\oplus(-\eta_R(b))}{%
    \FormulaRefAuto{IntegerRingPairValueDef}[def]{1,2,3}}
  \proofstepwidestarbreakkeep[1,4,5]{v_R(c,d)=\eta_R(c)\oplus(-\eta_R(d))}{%
    \FormulaRefAuto{IntegerRingPairValueDef}[def]{1,4,5}}
  \proofstepwidestarbreakkeep[1,2,3,4,5]{(\eta_R(a)\oplus(-\eta_R(b)))\odot(\eta_R(c)\oplus(-\eta_R(d)))=((\eta_R(a)\odot\eta_R(c))\oplus(\eta_R(b)\odot\eta_R(d)))\oplus(-((\eta_R(a)\odot\eta_R(d))\oplus(\eta_R(b)\odot\eta_R(c))))}{%
    \FormulaRefAuto{RingDifferenceMultiplication}[thm]{1,8,9,10,11}}
  \proofstepwidestarbreakkeep[1,2,3,4,5]{v_R(ac+bd,ad+bc)=v_R(a,b)\odot v_R(c,d)}{%
    \rIE{22,23,13,15,17,19,25,26,27,24}}
  \closeproofpartwide
\end{tabproofsplitwide}
\endgroup


\FormulaThmDeltaK[Die kanonische Ganzzahlabbildung erhält die Addition]{%
  \begin{aligned}[t]
    &\RingAlg{R}{\oplus}{\odot}{0_R}{1_R}\dsep z,w\in\mathbb Z\\
    &\vdash
      \IntegerRingMap{R,\oplus,\odot,0_R,1_R}(z+w)
      =
      \IntegerRingMap{R,\oplus,\odot,0_R,1_R}(z)
      \oplus
      \IntegerRingMap{R,\oplus,\odot,0_R,1_R}(w)
  \end{aligned}
}{IntegerRingMapAddition}{%
  \DeltaRow{Mengen}{R\dsep z\dsep w\dsep a\dsep b\dsep c\dsep d}
  \DeltaRow{Funktionen}{\oplus\dsep\odot}
}

\[
  \eta_R=\NaturalSemiringMap{R,\oplus,\odot,0_R,1_R},\qquad
  \Phi_R=\IntegerRingMap{R,\oplus,\odot,0_R,1_R},
  \qquad
  v_R(a,b)=\IntegerRingPairValue{R,\oplus,\odot,0_R,1_R}{a}{b}.
\]
\begingroup
\ProofWideReasonsNearFormula
\begin{tabproofsplitwide}
  \proofpartwide{Vertreter der beiden ganzen Zahlen}
  \proofstepwidestarbreakkeep[1]{\RingAlg{R}{\oplus}{\odot}{0_R}{1_R}}{%
    \rA}
  \proofstepwidestarbreakkeep[2]{z\in\mathbb Z}{%
    \rA}
  \proofstepwidestarbreakkeep[3]{w\in\mathbb Z}{%
    \rA}
  \proofstepwidestarbreakkeep[2]{\exists a,b\in\mathbb N\,(z=\IntClass{a}{b})}{%
    \FormulaRefAuto{IntegerPairRepresentative}[thm]{2}}
  \proofstepwidestarbreakkeep[3]{\exists c,d\in\mathbb N\,(w=\IntClass{c}{d})}{%
    \FormulaRefAuto{IntegerPairRepresentative}[thm]{3}}
  \proofstepwidestarbreakkeep[6]{a,b\in\mathbb N\land z=\IntClass{a}{b}}{%
    \rA}
  \proofstepwidestarbreakkeep[7]{c,d\in\mathbb N\land w=\IntClass{c}{d}}{%
    \rA}
  \proofstepwidestarbreakkeep[6]{a\in\mathbb N}{%
    \rAEa{\rAEa{6}}}
  \proofstepwidestarbreakkeep[6]{b\in\mathbb N}{%
    \rAEb{\rAEa{6}}}
  \proofstepwidestarbreakkeep[7]{c\in\mathbb N}{%
    \rAEa{\rAEa{7}}}
  \proofstepwidestarbreakkeep[7]{d\in\mathbb N}{%
    \rAEb{\rAEa{7}}}
  \proofstepwidestarbreakkeep[6,7]{a+c\in\mathbb N}{%
    \FormulaRefAuto{PeanoAddClosure}[thm]{8,10}}
  \proofstepwidestarbreakkeep[6,7]{b+d\in\mathbb N}{%
    \FormulaRefAuto{PeanoAddClosure}[thm]{9,11}}
  \closeproofpartwide

  \proofpartwide{Auswertung und Operationsgesetz der Paarwerte}
  \setcounter{proofstepnr}{13}
  \proofstepwidestarbreakkeep[6,7]{z+w=\IntClass{a+c}{b+d}}{%
    \FormulaRefAuto{IntegerAdditionDef}[def]{6,7}}
  \proofstepwidestarbreakkeep[1,6,7]{\Phi_R(\IntClass{a+c}{b+d})=v_R(a+c,b+d)}{%
    \FormulaRefAuto{IntegerRingMapDef}[def]{1,12,13}}
  \proofstepwidestarbreakkeep[1,6,7]{\Phi_R(z+w)=v_R(a+c,b+d)}{%
    \rIE{14,15}}
  \proofstepwidestarbreakkeep[1,6,7]{v_R(a+c,b+d)=v_R(a,b)\oplus v_R(c,d)}{%
    \FormulaRefAuto{IntegerRingPairValueAddition}[thm]{1,8,9,10,11}}
  \proofstepwidestarbreakkeep[1,6]{\Phi_R(\IntClass{a}{b})=v_R(a,b)}{%
    \FormulaRefAuto{IntegerRingMapDef}[def]{1,8,9}}
  \proofstepwidestarbreakkeep[1,7]{\Phi_R(\IntClass{c}{d})=v_R(c,d)}{%
    \FormulaRefAuto{IntegerRingMapDef}[def]{1,10,11}}
  \proofstepwidestarbreakkeep[1,6]{\Phi_R(z)=v_R(a,b)}{%
    \rIE{\rAEb{6},18}}
  \proofstepwidestarbreakkeep[1,7]{\Phi_R(w)=v_R(c,d)}{%
    \rIE{\rAEb{7},19}}
  \proofstepwidestarbreakkeep[1,6,7]{\Phi_R(z+w)=\Phi_R(z)\oplus\Phi_R(w)}{%
    \rIE{20,21,\rChain{16,17}}}
  \closeproofpartwide

  \proofpartwide{Existenzannahmen schließen}
  \setcounter{proofstepnr}{22}
  \proofstepwidestarbreakkeep[1,3,6]{\Phi_R(z+w)=\Phi_R(z)\oplus\Phi_R(w)}{%
    \rEEStar{5,7,22}}
  \proofstepwidestarbreakkeep[1,2,3]{\Phi_R(z+w)=\Phi_R(z)\oplus\Phi_R(w)}{%
    \rEEStar{4,6,23}}
  \closeproofpartwide
\end{tabproofsplitwide}
\endgroup


\FormulaThmDeltaK[Die kanonische Ganzzahlabbildung erhält die Multiplikation]{%
  \begin{aligned}[t]
    &\RingAlg{R}{\oplus}{\odot}{0_R}{1_R}\dsep z,w\in\mathbb Z\\
    &\vdash
      \IntegerRingMap{R,\oplus,\odot,0_R,1_R}(z\cdot w)
      =
      \IntegerRingMap{R,\oplus,\odot,0_R,1_R}(z)
      \odot
      \IntegerRingMap{R,\oplus,\odot,0_R,1_R}(w)
  \end{aligned}
}{IntegerRingMapMultiplication}{%
  \DeltaRow{Mengen}{R\dsep z\dsep w\dsep a\dsep b\dsep c\dsep d}
  \DeltaRow{Funktionen}{\oplus\dsep\odot}
}

\[
  \eta_R=\NaturalSemiringMap{R,\oplus,\odot,0_R,1_R},\qquad
  \Phi_R=\IntegerRingMap{R,\oplus,\odot,0_R,1_R},
  \qquad
  v_R(a,b)=\IntegerRingPairValue{R,\oplus,\odot,0_R,1_R}{a}{b}.
\]
\begingroup
\ProofWideReasonsNearFormula
\begin{tabproofsplitwide}
  \proofpartwide{Vertreter der beiden ganzen Zahlen}
  \proofstepwidestarbreakkeep[1]{\RingAlg{R}{\oplus}{\odot}{0_R}{1_R}}{%
    \rA}
  \proofstepwidestarbreakkeep[2]{z\in\mathbb Z}{%
    \rA}
  \proofstepwidestarbreakkeep[3]{w\in\mathbb Z}{%
    \rA}
  \proofstepwidestarbreakkeep[2]{\exists a,b\in\mathbb N\,(z=\IntClass{a}{b})}{%
    \FormulaRefAuto{IntegerPairRepresentative}[thm]{2}}
  \proofstepwidestarbreakkeep[3]{\exists c,d\in\mathbb N\,(w=\IntClass{c}{d})}{%
    \FormulaRefAuto{IntegerPairRepresentative}[thm]{3}}
  \proofstepwidestarbreakkeep[6]{a,b\in\mathbb N\land z=\IntClass{a}{b}}{%
    \rA}
  \proofstepwidestarbreakkeep[7]{c,d\in\mathbb N\land w=\IntClass{c}{d}}{%
    \rA}
  \proofstepwidestarbreakkeep[6]{a\in\mathbb N}{%
    \rAEa{\rAEa{6}}}
  \proofstepwidestarbreakkeep[6]{b\in\mathbb N}{%
    \rAEb{\rAEa{6}}}
  \proofstepwidestarbreakkeep[7]{c\in\mathbb N}{%
    \rAEa{\rAEa{7}}}
  \proofstepwidestarbreakkeep[7]{d\in\mathbb N}{%
    \rAEb{\rAEa{7}}}
  \proofstepwidestarbreakkeep[6,7]{ac\in\mathbb N}{%
    \FormulaRefAuto{PeanoMulClosure}[thm]{8,10}}
  \proofstepwidestarbreakkeep[6,7]{bd\in\mathbb N}{%
    \FormulaRefAuto{PeanoMulClosure}[thm]{9,11}}
  \proofstepwidestarbreakkeep[6,7]{ad\in\mathbb N}{%
    \FormulaRefAuto{PeanoMulClosure}[thm]{8,11}}
  \proofstepwidestarbreakkeep[6,7]{bc\in\mathbb N}{%
    \FormulaRefAuto{PeanoMulClosure}[thm]{9,10}}
  \proofstepwidestarbreakkeep[6,7]{ac+bd\in\mathbb N}{%
    \FormulaRefAuto{PeanoAddClosure}[thm]{12,13}}
  \proofstepwidestarbreakkeep[6,7]{ad+bc\in\mathbb N}{%
    \FormulaRefAuto{PeanoAddClosure}[thm]{14,15}}
  \closeproofpartwide

  \proofpartwide{Auswertung und Operationsgesetz der Paarwerte}
  \setcounter{proofstepnr}{17}
  \proofstepwidestarbreakkeep[6,7]{z\cdot w=\IntClass{ac+bd}{ad+bc}}{%
    \FormulaRefAuto{IntegerMultiplicationDef}[def]{6,7}}
  \proofstepwidestarbreakkeep[1,6,7]{\Phi_R(\IntClass{ac+bd}{ad+bc})=v_R(ac+bd,ad+bc)}{%
    \FormulaRefAuto{IntegerRingMapDef}[def]{1,16,17}}
  \proofstepwidestarbreakkeep[1,6,7]{\Phi_R(z\cdot w)=v_R(ac+bd,ad+bc)}{%
    \rIE{18,19}}
  \proofstepwidestarbreakkeep[1,6,7]{v_R(ac+bd,ad+bc)=v_R(a,b)\odot v_R(c,d)}{%
    \FormulaRefAuto{IntegerRingPairValueMultiplication}[thm]{1,8,9,10,11}}
  \proofstepwidestarbreakkeep[1,6]{\Phi_R(\IntClass{a}{b})=v_R(a,b)}{%
    \FormulaRefAuto{IntegerRingMapDef}[def]{1,8,9}}
  \proofstepwidestarbreakkeep[1,7]{\Phi_R(\IntClass{c}{d})=v_R(c,d)}{%
    \FormulaRefAuto{IntegerRingMapDef}[def]{1,10,11}}
  \proofstepwidestarbreakkeep[1,6]{\Phi_R(z)=v_R(a,b)}{%
    \rIE{\rAEb{6},22}}
  \proofstepwidestarbreakkeep[1,7]{\Phi_R(w)=v_R(c,d)}{%
    \rIE{\rAEb{7},23}}
  \proofstepwidestarbreakkeep[1,6,7]{\Phi_R(z\cdot w)=\Phi_R(z)\odot\Phi_R(w)}{%
    \rIE{24,25,\rChain{20,21}}}
  \closeproofpartwide

  \proofpartwide{Existenzannahmen schließen}
  \setcounter{proofstepnr}{26}
  \proofstepwidestarbreakkeep[1,3,6]{\Phi_R(z\cdot w)=\Phi_R(z)\odot\Phi_R(w)}{%
    \rEEStar{5,7,26}}
  \proofstepwidestarbreakkeep[1,2,3]{\Phi_R(z\cdot w)=\Phi_R(z)\odot\Phi_R(w)}{%
    \rEEStar{4,6,27}}
  \closeproofpartwide
\end{tabproofsplitwide}
\endgroup


\FormulaThmDeltaK[Null- und Einserhaltung der kanonischen Ganzzahlabbildung]{%
  \RingAlg{R}{\oplus}{\odot}{0_R}{1_R}
  \vdash
  \IntegerRingMap{R,\oplus,\odot,0_R,1_R}(\IntZero)=0_R
  \land
  \IntegerRingMap{R,\oplus,\odot,0_R,1_R}(\IntOne)=1_R%
}{IntegerRingMapConstants}{%
  \DeltaRow{Mengen}{R\dsep\IntZero\dsep\IntOne}
  \DeltaRow{Funktionen}{\oplus\dsep\odot}
}

\[
  \eta_R=\NaturalSemiringMap{R,\oplus,\odot,0_R,1_R},\qquad
  \Phi_R=\IntegerRingMap{R,\oplus,\odot,0_R,1_R},
  \qquad
  v_R(a,b)=\IntegerRingPairValue{R,\oplus,\odot,0_R,1_R}{a}{b}.
\]
\begingroup
\ProofWideReasonsNearFormula
\begin{tabproofsplitwide}
  \proofpartwide{Träger und natürliche Konstanten}
  \proofstepwidestarbreakkeep[1]{\RingAlg{R}{\oplus}{\odot}{0_R}{1_R}}{%
    \rA}
  \proofstepwidestarbreakkeep[]{0\in\mathbb N}{%
    \FormulaRefAuto{PeanoZero}[ax]}
  \proofstepwidestarbreakkeep[]{1\in\mathbb N}{%
    \FormulaRefAuto{PeanoOneInN}[thm]}
  \proofstepwidestarbreakkeep[1]{\SemiringAlg{R}{\oplus}{\odot}{0_R}{1_R}}{%
    \FormulaRefAuto{RingUnderlyingSemiring}[thm]{1}}
  \proofstepwidestarbreakkeep[1]{\GroupAlg{R}{\oplus}{0_R}}{%
    \FormulaRefAuto{RingAdditiveGroup}[thm]{1}}
  \proofstepwidestarbreakkeep[1]{0_R\in R}{%
    \FormulaRefAuto{GroupIdentityCarrier}[thm]{5}}
  \proofstepwidestarbreakkeep[1]{\eta_R(0)=0_R}{%
    \FormulaRefAuto{NaturalSemiringMapZero}[thm]{4}}
  \proofstepwidestarbreakkeep[1]{\eta_R(1)=1_R}{%
    \FormulaRefAuto{NaturalSemiringMapOne}[thm]{4}}
  \closeproofpartwide

  \proofpartwide{Nullerhaltung}
  \setcounter{proofstepnr}{8}
  \proofstepwidestarbreakkeep[]{\IntZero=\IntClass{0}{0}}{%
    \FormulaRefAuto{IntegerZeroClass}[thm]}
  \proofstepwidestarbreakkeep[1]{\Phi_R(\IntClass{0}{0})=\eta_R(0)\oplus(-\eta_R(0))}{%
    \FormulaRefAuto{IntegerRingMapClassValue}[thm]{1,2,2}}
  \proofstepwidestarbreakkeep[1]{0_R\oplus(-0_R)=0_R}{%
    \FormulaRefAuto{RingNegativeRight}[thm]{1,6}}
  \proofstepwidestarbreakkeep[1]{\Phi_R(\IntZero)=0_R}{%
    \rIE{9,7,11,10}}
  \closeproofpartwide

  \proofpartwide{Einserhaltung}
  \setcounter{proofstepnr}{12}
  \proofstepwidestarbreakkeep[1]{-0_R\in R}{%
    \FormulaRefAuto{RingNegativeCarrier}[thm]{1,6}}
  \proofstepwidestarbreakkeep[1]{0_R\oplus(-0_R)=-0_R}{%
    \FormulaRefAuto{GroupLeftIdentity}[thm]{5,13}}
  \proofstepwidestarbreakkeep[1]{-0_R=0_R}{%
    \FormulaRefAuto{a=b\dsep a=c\vdash b=c}[thm]{14,11}}
  \proofstepwidestarbreakkeep[]{\IntOne=\IntEmbed(1)}{%
    \FormulaRefAuto{IntegerOneDef}[def]}
  \proofstepwidestarbreakkeep[]{\IntEmbed(1)=\IntClass{1}{0}}{%
    \FormulaRefAuto{NaturalIntegerEmbeddingValue}[thm]{3}}
  \proofstepwidestarbreakkeep[]{\IntOne=\IntClass{1}{0}}{%
    \rChain{16,17}}
  \proofstepwidestarbreakkeep[1]{\Phi_R(\IntClass{1}{0})=\eta_R(1)\oplus(-\eta_R(0))}{%
    \FormulaRefAuto{IntegerRingMapClassValue}[thm]{1,3,2}}
  \proofstepwidestarbreakkeep[1]{\MonoidAlg{R}{\odot}{1_R}}{%
    \FormulaRefAuto{RingMultiplicativeMonoidAxiom}[ax]{1}}
  \proofstepwidestarbreakkeep[1]{1_R\in R}{%
    \FormulaRefAuto{MonoidIdentityCarrierAxiom}[ax]{20}}
  \proofstepwidestarbreakkeep[1]{1_R\oplus0_R=1_R}{%
    \FormulaRefAuto{GroupRightIdentity}[thm]{5,21}}
  \proofstepwidestarbreakkeep[1]{\Phi_R(\IntOne)=1_R}{%
    \rIE{18,8,7,15,22,19}}
  \proofstepwidestarbreakkeep[1]{\Phi_R(\IntZero)=0_R\land\Phi_R(\IntOne)=1_R}{%
    \rAI{12,23}}
  \closeproofpartwide
\end{tabproofsplitwide}
\endgroup


\FormulaThmDeltaK[Die kanonische Ganzzahlabbildung ist ein Ringhomomorphismus]{%
  \RingAlg{R}{\oplus}{\odot}{0_R}{1_R}
  \vdash
  \RingHom{\IntegerRingMap{R,\oplus,\odot,0_R,1_R}}
    {\mathbb Z,+,\cdot,\IntZero,\IntOne}
    {R,\oplus,\odot,0_R,1_R}%
}{IntegerRingCanonicalHomomorphism}{%
  \DeltaRow{Mengen}{R\dsep z\dsep w}
  \DeltaRow{Funktionen}{\oplus\dsep\odot}
}

Für die beiden folgenden Beweise setzen wir lokal
\[
  \mathfrak Z
    \coloneqq(\mathbb Z,+,\cdot,\IntZero,\IntOne),
  \qquad
  \mathfrak R
    \coloneqq(R,\oplus,\odot,0_R,1_R),
  \qquad
  \Phi_R
    \coloneqq\IntegerRingMap{R,\oplus,\odot,0_R,1_R}.
\]

\begin{tabproofsplitwide}
  \proofpartwide{Additiver Gruppenhomomorphismus}
  \proofstepwidestar[1]{%
    \RingAlg{R}{\oplus}{\odot}{0_R}{1_R}}{\rA}
  \proofstepwidestar[]{%
    \RingAlg{\mathbb Z}{+}{\cdot}{\IntZero}{\IntOne}}{%
    \FormulaRefAuto{IntegerUnitalRing}[thm]}
  \proofstepwidestar[1]{%
    \Phi_R\colon\mathbb Z\to R}{%
    \FormulaRefAuto{IntegerRingMapFunction}[thm]{1}}
  \proofstepwidestar[]{%
    \GroupAlg{\mathbb Z}{+}{\IntZero}}{%
    \FormulaRefAuto{IntegerAdditiveGroup}[thm]}
  \proofstepwidestar[1]{\GroupAlg{R}{\oplus}{0_R}}{%
    \FormulaRefAuto{RingAdditiveGroup}[thm]{1}}
  \proofstepwidestar[]{\SemigroupAlg{\mathbb Z}{+}}{%
    \FormulaRefAuto{GroupBaseSemigroup}[thm]{4}}
  \proofstepwidestar[1]{\SemigroupAlg{R}{\oplus}}{%
    \FormulaRefAuto{GroupBaseSemigroup}[thm]{5}}
  \proofstepwidestar[8]{z\in\mathbb Z}{\rA}
  \proofstepwidestar[9]{w\in\mathbb Z}{\rA}
  \proofstepwidestar[1,8,9]{%
    \begin{aligned}[t]
      &\Phi_R(z+w)\\[-2pt]
      &\quad=
      \Phi_R(z)
      \oplus
      \Phi_R(w)
    \end{aligned}}{%
    \FormulaRefAuto{IntegerRingMapAddition}[thm]{1,8,9}}
  \proofstepwidestar[1]{%
    \SemigroupHom{\Phi_R}
      {\mathbb Z,+}{R,\oplus}}{%
    \FormulaRefAuto{SemigroupHomDef}[def]{6,7,3,10}}
  \proofstepwidestar[1]{%
    \GroupHom{\Phi_R}
      {\mathbb Z,+,\IntZero}{R,\oplus,0_R}}{%
    \FormulaRefAuto{GroupHomDef}[def]{4,5,11}}
  \closeproofpartwide

  \proofpartwide{Multiplikativer Monoidanteil und Ringhomomorphismus}
  \setcounter{proofstepnr}{12}
  \proofstepwidestar[]{%
    \MonoidAlg{\mathbb Z}{\cdot}{\IntOne}}{%
    \FormulaRefAuto{IntegerMultiplicativeMonoid}[thm]}
  \proofstepwidestar[1]{\MonoidAlg{R}{\odot}{1_R}}{%
    \FormulaRefAuto{RingMultiplicativeMonoidAxiom}[ax]{1}}
  \proofstepwidestar[]{\SemigroupAlg{\mathbb Z}{\cdot}}{%
    \FormulaRefAuto{MonoidBaseSemigroupAxiom}[ax]{13}}
  \proofstepwidestar[1]{\SemigroupAlg{R}{\odot}}{%
    \FormulaRefAuto{MonoidBaseSemigroupAxiom}[ax]{14}}
  \proofstepwidestar[1,8,9]{%
    \begin{aligned}[t]
      &\Phi_R(z\cdot w)\\[-2pt]
      &\quad=
      \Phi_R(z)
      \odot
      \Phi_R(w)
    \end{aligned}}{%
    \FormulaRefAuto{IntegerRingMapMultiplication}[thm]{1,8,9}}
  \proofstepwidestar[1]{%
    \SemigroupHom{\Phi_R}
      {\mathbb Z,\cdot}{R,\odot}}{%
    \FormulaRefAuto{SemigroupHomDef}[def]{15,16,3,17}}
  \proofstepwidestar[1]{%
    \Phi_R(\IntZero)=0_R
    \land
    \Phi_R(\IntOne)=1_R}{%
    \FormulaRefAuto{IntegerRingMapConstants}[thm]{1}}
  \proofstepwidestar[1]{%
    \Phi_R(\IntOne)=1_R}{%
    \rAEb{19}}
  \proofstepwidestar[1]{%
    \MonoidHom{\Phi_R}
      {\mathbb Z,\cdot,\IntOne}{R,\odot,1_R}}{%
    \FormulaRefAuto{MonoidHomDef}[def]{13,14,18,20}}
  \proofstepwidestar[1]{%
    \RingHom{\Phi_R}{\mathfrak Z}{\mathfrak R}}{%
    \FormulaRefAuto{RingHomDef}[def]{2,1,12,21}}
  \closeproofpartwide
\end{tabproofsplitwide}

\FormulaThmDeltaK[Universelle Eigenschaft der ganzen Zahlen]{%
  \begin{aligned}[t]
    &\RingAlg{R}{\oplus}{\odot}{0_R}{1_R}\vdash{}\\[-2pt]
    &\exists! f\,
      \RingHom{f}
        {\mathbb Z,+,\cdot,\IntZero,\IntOne}
        {R,\oplus,\odot,0_R,1_R}
  \end{aligned}
}{IntegerRingUniversalHomomorphism}{%
  \DeltaRow{Mengen}{R\dsep z\dsep m\dsep n}
  \DeltaRow{Funktionen}{f\dsep\oplus\dsep\odot}
}

\begin{tabproofsplitwide}
  \proofpartwideindR[Eindeutigkeit auf der natürlichen Einbettung]{%
    \begin{aligned}[t]
      &\operatorname{Ring}(\mathfrak R)\dsep\\[-2pt]
      &\RingHom{f}{\mathfrak Z}{\mathfrak R}\dsep\\[-2pt]
      &n\in\mathbb N\vdash
        f(\IntEmbed(n))={}\\[-2pt]
      &\quad
        \NaturalSemiringMap{R,\oplus,\odot,0_R,1_R}(n)
    \end{aligned}
  }{%
    \begin{aligned}[t]
      &\operatorname{Ring}(\mathfrak R)\dsep\\[-2pt]
      &\RingHom{f}{\mathfrak Z}{\mathfrak R}\dsep\\[-2pt]
      &n\in\mathbb N\vdash
        f(\IntEmbed(n))={}\\[-2pt]
      &\quad
        \NaturalSemiringMap{R,\oplus,\odot,0_R,1_R}(n)
    \end{aligned}
  }[IntegerRingHomEmbeddingValue]
    \proofstepwidestar[1]{%
      \operatorname{Ring}(\mathfrak R)}{\rA}
    \proofstepwidestar[2]{%
      \RingHom{f}{\mathfrak Z}{\mathfrak R}}{\rA}
    \proofstepwidestar[2]{f(\IntZero)=0_R}{%
      \FormulaRefAuto{RingHomZeroPreservation}[thm]{2}}
    \proofstepwidestar[]{\IntZero=\IntEmbed(0)}{%
      \FormulaRefAuto{IntegerZeroDef}[def]}
    \proofstepwidestar[2]{f(\IntEmbed(0))=0_R}{\rIE{4,3}}
    \proofstepwidestar[1]{\SemiringAlg{R}{\oplus}{\odot}{0_R}{1_R}}{%
      \FormulaRefAuto{RingUnderlyingSemiring}[thm]{1}}
    \proofstepwidestar[1]{%
      \NaturalSemiringMap{R,\oplus,\odot,0_R,1_R}(0)=0_R}{%
      \FormulaRefAuto{NaturalSemiringMapZero}[thm]{%
        6}}
    \proofstepwidestar[1,2]{%
      \begin{aligned}[t]
        &f(\IntEmbed(0))={}\\[-2pt]
        &\quad\NaturalSemiringMap{R,\oplus,\odot,0_R,1_R}(0)
      \end{aligned}}{%
      \rIE{5,\FormulaRefAuto{a=b\vdash b=a}{7}}}
    \proofstepwidestar[9]{m\in\mathbb N}{\rA}
    \proofstepwidestar[10]{%
      \begin{aligned}[t]
        &f(\IntEmbed(m))={}\\[-2pt]
        &\quad\NaturalSemiringMap{R,\oplus,\odot,0_R,1_R}(m)
      \end{aligned}}{\rA}
    \proofstepwidestar[]{1\in\mathbb N}{%
      \FormulaRefAuto{PeanoOneInN}[thm]}
    \proofstepwidestar[9]{%
      \IntEmbed(m+1)=\IntEmbed(m)+\IntEmbed(1)}{%
      \BandXXIXProofRefStack{%
        \FormulaRefAuto{NaturalIntegerEmbeddingAdditive}[thm]\\
        \bigl(9,11\bigr)}}
    \proofstepwidestar[9]{%
      \begin{aligned}[t]
        &f(\IntEmbed(m+1))={}\\[-2pt]
        &\quad f\bigl(\IntEmbed(m)+\IntEmbed(1)\bigr)
      \end{aligned}}{\rIE{12,\rII}}
    \proofstepwidestar[]{1\in\mathbb N}{%
      \FormulaRefAuto{PeanoOneInN}[thm]}
    \proofstepwidestar[]{\IntEmbed\colon\mathbb N\to\mathbb Z}{%
      \FormulaRefAuto{NaturalIntegerEmbeddingFunction}[thm]}
    \proofstepwidestar[]{\IntEmbed\colon\mathbb N\to\mathbb Z}{%
      \FormulaRefAuto{NaturalIntegerEmbeddingFunction}[thm]}
    \proofstepwidestar[2,9]{\IntEmbed(m)\in\mathbb Z}{%
      \FormulaRefAuto{F\colon A\to B\dsep x\in A\vdash F(x)\in B}{%
              15,9}}
    \proofstepwidestar[2,9]{\IntEmbed(1)\in\mathbb Z}{%
      \FormulaRefAuto{F\colon A\to B\dsep x\in A\vdash F(x)\in B}{%
              16,
              14}}
    \proofstepwidestar[2,9]{%
      \begin{aligned}[t]
        &f\bigl(\IntEmbed(m)+\IntEmbed(1)\bigr)={}\\[-2pt]
        &\quad f(\IntEmbed(m))\oplus f(\IntEmbed(1))
      \end{aligned}}{%
      \BandXXIXProofRefStack{%
        \FormulaRefAuto{RingHomAdditionAxiom}[ax]\\
        \left(
          \begin{gathered}[t]
            2,\,
            17,\\[-2pt]
            18
          \end{gathered}
        \right)}}
    \proofstepwidestar[2]{\IntOne=\IntEmbed(1)}{%
      \FormulaRefAuto{IntegerOneDef}[def]}
    \proofstepwidestar[2]{f(\IntOne)=1_R}{%
      \FormulaRefAuto{RingHomOneAxiom}[ax]{2}}
    \proofstepwidestar[2]{f(\IntEmbed(1))=1_R}{%
      \rIE{%
        20,
        21}}
    \proofstepwidestar[1,2,10]{%
      \begin{aligned}[t]
        &f(\IntEmbed(m))\oplus f(\IntEmbed(1))={}\\[-2pt]
        &\quad\NaturalSemiringMap{R,\oplus,\odot,0_R,1_R}(m)\oplus1_R
      \end{aligned}}{%
      \rIE{10,22}}
    \proofstepwidestar[1]{\SemiringAlg{R}{\oplus}{\odot}{0_R}{1_R}}{%
      \FormulaRefAuto{RingUnderlyingSemiring}[thm]{1}}
    \proofstepwidestar[1,9]{\begin{aligned}[t]&\NaturalSemiringMap{R,\oplus,\odot,0_R,1_R}(m+1)={}\\[-2pt]&\quad\NaturalSemiringMap{R,\oplus,\odot,0_R,1_R}(m)\oplus1_R\end{aligned}}{%
      \FormulaRefAuto{NaturalSemiringMapAddOne}[thm]{24,9}}
    \proofstepwidestar[1,9]{%
      \begin{aligned}[t]
        &\NaturalSemiringMap{R,\oplus,\odot,0_R,1_R}(m)\oplus1_R={}\\[-2pt]
        &\quad\NaturalSemiringMap{R,\oplus,\odot,0_R,1_R}(m+1)
      \end{aligned}}{\FormulaRefAuto{a=b\vdash b=a}[thm]{25}}
    \proofstepwidestar[1,2,9,10]{%
      \begin{aligned}[t]
        &f(\IntEmbed(m+1))={}\\[-2pt]
        &\quad\NaturalSemiringMap{R,\oplus,\odot,0_R,1_R}(m+1)
      \end{aligned}}{%
      \rChain{13,19,23,26}}
    \proofstepwidestar[28]{n\in\mathbb N}{\rA}
    \proofstepwidestar[1,2,28]{%
      \begin{aligned}[t]
        &f(\IntEmbed(n))={}\\[-2pt]
        &\quad\NaturalSemiringMap{R,\oplus,\odot,0_R,1_R}(n)
      \end{aligned}}{%
      \FormulaRefAuto{PeanoInductionRuleAddOne}[thm]{8,27,28}}
  \closeproofpartwideindR

  \proofpartwideindR[Punktweise Eindeutigkeit auf den ganzen Zahlen]{%
    \begin{aligned}[t]
      &\operatorname{Ring}(\mathfrak R)\dsep\\[-2pt]
      &\RingHom{f}{\mathfrak Z}{\mathfrak R}\dsep\\[-2pt]
      &z\in\mathbb Z\vdash f(z)=\Phi_R(z)
    \end{aligned}
  }{%
    \begin{aligned}[t]
      &\operatorname{Ring}(\mathfrak R)\dsep\\[-2pt]
      &\RingHom{f}{\mathfrak Z}{\mathfrak R}\dsep\\[-2pt]
      &z\in\mathbb Z\vdash f(z)=\Phi_R(z)
    \end{aligned}
  }[IntegerRingHomPointwiseUnique]
    \proofstepwidestar[1]{%
      \operatorname{Ring}(\mathfrak R)}{\rA}
    \proofstepwidestar[2]{%
      \RingHom{f}{\mathfrak Z}{\mathfrak R}}{\rA}
    \proofstepwidestar[3]{z\in\mathbb Z}{\rA}
    \proofstepwidestar[1]{%
      \RingHom{\Phi_R}{\mathfrak Z}{\mathfrak R}}{%
      \FormulaRefAuto{IntegerRingCanonicalHomomorphism}[thm]{1}}
    \proofstepwidestar[3]{%
      \exists n\in\mathbb N\,
      \bigl(z=\IntEmbed(n)\lor z=-\IntEmbed(n)\bigr)}{%
      \FormulaRefAuto{IntegerSignedNormalForm}[thm]{3}}
    \proofstepwidestar[6]{%
      n\in\mathbb N\land
      \bigl(z=\IntEmbed(n)\lor z=-\IntEmbed(n)\bigr)}{\rA}
    \proofstepwidestar[6]{n\in\mathbb N}{\rAEa{6}}
    \proofstepwidestar[6]{%
      z=\IntEmbed(n)\lor z=-\IntEmbed(n)}{\rAEb{6}}
    \proofstepwidestar[1,2,6]{%
      \begin{aligned}[t]
        &f(\IntEmbed(n))={}\\[-2pt]
        &\quad\NaturalSemiringMap{R,\oplus,\odot,0_R,1_R}(n)
      \end{aligned}}{%
      \FormulaRefAuto{IntegerRingHomEmbeddingValue}[thm]{1,2,7}}
    \proofstepwidestar[1,6]{%
      \begin{aligned}[t]
        &\Phi_R(\IntEmbed(n))={}\\[-2pt]
        &\quad\NaturalSemiringMap{R,\oplus,\odot,0_R,1_R}(n)
      \end{aligned}}{%
      \FormulaRefAuto{IntegerRingHomEmbeddingValue}[thm]{1,4,7}}
    \proofstepwidestar[1,2,6]{%
      f(\IntEmbed(n))=\Phi_R(\IntEmbed(n))}{%
      \FormulaRefAuto{a=b\dsep c=b\vdash a=c}{9,10}}
    \proofstepwidestar[12]{z=\IntEmbed(n)}{\rA}
    \proofstepwidestar[1,2,6,12]{%
      f(z)=\Phi_R(z)}{%
      \rIE{12,11}}
    \proofstepwidestar[14]{z=-\IntEmbed(n)}{\rA}
    \proofstepwidestar[]{\IntEmbed\colon\mathbb N\to\mathbb Z}{%
      \FormulaRefAuto{NaturalIntegerEmbeddingFunction}[thm]}
    \proofstepwidestar[6]{\IntEmbed(n)\in\mathbb Z}{%
      \BandXXIXProofRefStack{%
        \FormulaRefAuto{F\colon A\to B\dsep x\in A\vdash F(x)\in B}[thm]\\
        \bigl(15,7\bigr)}}
    \proofstepwidestar[2,6]{%
      f(-\IntEmbed(n))=-f(\IntEmbed(n))}{%
      \FormulaRefAuto{RingHomNegativePreservation}[thm]{2,16}}
    \proofstepwidestar[1,6]{%
      \Phi_R(-\IntEmbed(n))=-\Phi_R(\IntEmbed(n))}{%
      \FormulaRefAuto{RingHomNegativePreservation}[thm]{4,16}}
    \proofstepwidestar[1,2,6]{%
      -f(\IntEmbed(n))=-\Phi_R(\IntEmbed(n))}{%
      \rIE{11,\rII}}
    \proofstepwidestar[1,6]{%
      -\Phi_R(\IntEmbed(n))=\Phi_R(-\IntEmbed(n))}{%
      \FormulaRefAuto{a=b\vdash b=a}{18}}
    \proofstepwidestar[1,2,6]{%
      f(-\IntEmbed(n))=\Phi_R(-\IntEmbed(n))}{%
      \rChain{17,19,20}}
    \proofstepwidestar[1,2,6,14]{%
      f(z)=\Phi_R(z)}{%
      \rIE{14,21}}
    \proofstepwidestar[1,2,6]{%
      f(z)=\Phi_R(z)}{%
      \rOE{8,12,13,14,22}}
    \proofstepwidestar[1,2,3]{%
      f(z)=\Phi_R(z)}{%
      \rEE{5,6,23}}
  \closeproofpartwideindR

  \proofpartwideindR[Eindeutigkeit der Ringabbildung]{%
    \begin{aligned}[t]
      &\operatorname{Ring}(\mathfrak R)\dsep\\[-2pt]
      &\RingHom{f}{\mathfrak Z}{\mathfrak R}\\
      &\vdash f=\Phi_R
    \end{aligned}
  }{%
    \begin{aligned}[t]
      &\operatorname{Ring}(\mathfrak R)\dsep\\[-2pt]
      &\RingHom{f}{\mathfrak Z}{\mathfrak R}\\
      &\vdash f=\Phi_R
    \end{aligned}
  }[IntegerRingHomUnique]
    \proofstepwidestar[1]{%
      \operatorname{Ring}(\mathfrak R)}{\rA}
    \proofstepwidestar[2]{%
      \RingHom{f}{\mathfrak Z}{\mathfrak R}}{\rA}
    \proofstepwidestar[2]{f\colon\mathbb Z\to R}{%
      \FormulaRefAuto{RingHomFunctionAxiom}[ax]{2}}
    \proofstepwidestar[1]{\Phi_R\colon\mathbb Z\to R}{%
      \FormulaRefAuto{IntegerRingMapFunction}[thm]{1}}
    \proofstepwidestar[5]{z\in\mathbb Z}{\rA}
    \proofstepwidestar[1,2,5]{%
      f(z)=\Phi_R(z)}{%
      \FormulaRefAuto{IntegerRingHomPointwiseUnique}[thm]{1,2,5}}
    \proofstepwidestar[1,2]{%
      \begin{aligned}[t]
        &\forall z\in\mathbb Z\,\\[-2pt]
        &f(z)=\Phi_R(z)
      \end{aligned}}{%
      \rUI{\rRI{5,6}}}
    \proofstepwidestar[1,2]{%
      f=\Phi_R}{%
      \FormulaRefAuto{FunctionExtensionality}[thm]{3,4,7}}
  \closeproofpartwideindR

  \proofpartwide{Existenz und Eindeutigkeit}
    \proofstepwidestar[1]{%
      \operatorname{Ring}(\mathfrak R)}{\rA}
    \proofstepwidestar[1]{%
      \RingHom{\Phi_R}{\mathfrak Z}{\mathfrak R}}{%
      \FormulaRefAuto{IntegerRingCanonicalHomomorphism}[thm]{1}}
    \proofstepwidestar[3]{%
      \RingHom{f}{\mathfrak Z}{\mathfrak R}}{\rA}
    \proofstepwidestar[1,3]{%
      f=\Phi_R}{%
      \FormulaRefAuto{IntegerRingHomUnique}[thm]{1,3}}
    \proofstepwidestar[1]{%
      \begin{aligned}[t]
        &\exists! f\,\\[-2pt]
        &\RingHom{f}{\mathfrak Z}{\mathfrak R}
      \end{aligned}}{\UEI{2,3,4}}
  \closeproofpartwide
\end{tabproofsplitwide}

\end{document}
